%% 第13回編集委員会（物理第1巻）  p1-07.tex  力-7　転がる・ロボットの腕：平面運動と機構（★4〜5・3段）　組 P1-4
%% 計画書 tools/phys1-plan.md の B「力-7」と A97（転がり）・A74（遠心力）を受けもつ。例題の数値は Python で確かめた
%%   （scratchpad/phys1/力-7/calc.py・calc2.py・calc3.py・calc4.py・ik_line.py）。g＝9.8 m/s²。
%% 出典（まだわかっていないこと・実機の値）
%%   ・自転車の自己安定：J. D. G. Kooijman, J. P. Meijaard, J. M. Papadopoulos, A. Ruina, A. L. Schwab,
%%     "A bicycle can be self-stable without gyroscopic or caster effects", Science 332, 339–342 (2011)
%%   ・ふろのうず：A. H. Shapiro, "Bath-tub vortex", Nature 196, 1080–1081 (1962)
%%     （MIT。直径約1.8 m・深さ約15 cmの丸い水そうで、水を24時間静かにしてから栓をぬくと左回り。MIT Technology Review の紹介記事で確認）
%%   ・産業用ロボットの位置くり返し精度：オムロン 産業用ロボット カタログ（水平多関節 Cobra 450 ±0.02 mm、
%%     eCobra ±0.017 mm、垂直多関節 Viper ±0.02〜0.03 mm）https://www.fa.omron.co.jp/data_pdf/selection/func_spec/industrial_robots_sbce-083_7_6_p6-7.pdf
%%   ・地球の自転の角速度：1恒星日 86164 s → 2π/86164 ＝ 7.29×10⁻⁵ rad/s（理科年表の恒星日）
%%   ・フーコーの振り子：1851年、パリ（パンテオンでの公開実験）。ふれる面が1周する時間＝恒星日÷sin（緯度）
%% 同じ題材：洗濯機の脱水（力-1 の例題：毎分800回転・半径 0.25 m・約 1800 m/s²）、カーブの摩擦 μ＝0.7（力-2 の例題）
\makeatletter
% 2リンクの腕  \pxharm{L1(mm)}{L2(mm)}{θ1}{θ2}{色}   点 (O)(E)(P) ができる
\newcommand\pxharm[5]{%
  \coordinate (O) at (0,0);
  \coordinate (E) at ({#1*cos(#3)},{#1*sin(#3)});
  \coordinate (P) at ($(E)+({#2*cos(#3+#4)},{#2*sin(#3+#4)})$);
  \draw[#5, line width=3pt, line cap=round] (O) -- (E) -- (P);
  \filldraw[fill=white, draw=bkInk, line width=0.7pt] (O) circle (1.5);
  \filldraw[fill=white, draw=bkInk, line width=0.7pt] (E) circle (1.3);
  \fill[bkInk] (P) circle (1.0);}
% 腕の土台  \pxhbase
\newcommand\pxhbase{\fill[bkGray!60] (-4.5,-3.2) rectangle (4.5,-1.4);\fill[bkGray!60] (-2,-1.6) rectangle (2,0);}
% 車輪  \pxhwheel{(x,y)}{半径}{スポークの角}
\newcommand\pxhwheel[3]{\begin{scope}[shift={#1}]
  \draw[bkInk!75, line width=1.5pt] (0,0) circle (#2);
  \foreach \k in {0,60,...,300} {\draw[bkGray, line width=0.5pt] (0,0) -- ({#3+\k}:{#2-0.8});}
  \fill[bkGray] (0,0) circle (0.9);
\end{scope}}
% 歯車  \pxhgear{(x,y)}{かみ合う円の半径(mm)}{歯の数}{回す角}{色}
\newcommand\pxhgear[5]{\begin{scope}[shift={#1}]
  \pgfmathsetmacro\pxhrr{#2-1.1}\pgfmathsetmacro\pxhro{#2+1.1}\pgfmathsetmacro\pxhw{90/#3}
  \fill[#5] (0,0) circle (\pxhrr);
  \foreach \k in {1,...,#3} {\pgfmathsetmacro\pxha{#4+\k*(360/#3)}
    \fill[#5] (\pxha-\pxhw*1.15:\pxhrr) -- (\pxha-\pxhw*0.55:\pxhro) -- (\pxha+\pxhw*0.55:\pxhro) -- (\pxha+\pxhw*1.15:\pxhrr) -- cycle;}
  \fill[white] (0,0) circle (1.5); \draw[bkInk, line width=0.5pt] (0,0) circle (1.5);
\end{scope}}
\makeatother

\phchapter{7}{転がる・ロボットの腕：\\平面運動と機構}{ロボットの手を、ねらった場所へ動かすには？}{\bklevelrange{4}{5}}{転がる・ロボットの腕：平面運動と機構}
\bkmichishirube{4-7,7-6,力-1,力-6}{力-16,波-3,波-4}

\begin{goal}
\begin{itemize}[label={\color{bkTheme}●}, leftmargin=1.6\zw, itemsep=0pt, topsep=1pt]
\item \bklevel{4}\enspace 転がる車輪の点の速さを「中心の速さ＋回る速さ」で求め、$v=r\omega$ を使える。
\item \bklevel{4}\enspace 2本の腕の関節の角から手の位置を、手の位置から関節の角を求められる。歯車の速さとトルクを求められる。
\item \bklevel{5}\enspace 平面運動の式と回る座標の式から、斜面を転がる加速度や、遠心力・コリオリの力を導ける。
\end{itemize}
\end{goal}

\phkoushiki{%
\phk{$v=r\omega$（すべらずに転がる車輪の中心の速さ）}{高校物理（発展）}{ph:力-7:korogari}
\phk{$a=\dfrac{g\sin\theta}{1+\frac{I}{mr^2}}$（斜面を転がる物の加速度）}{高校物理（発展）}{ph:力-7:shamen}
\phk{$mr\omega^2$（遠心力）}{高校物理}{ph:力-7:enshin}
\phkx{$v=r\omega$、$a=r\omega^2=\dfrac{v^2}{r}$（等速円運動）}{高校物理}{力-1 で導く}
\phkx{慣性力 $-m\vec A$}{高校物理}{力-2 で導く}
\phkx{$I\dfrac{\dd\omega}{\dd t}=N$、$\dfrac{1}{2}I\omega^2$}{高校物理（発展）}{力-6 で導く}
}

\begin{hatena}
\noindent\begin{minipage}[c]{0.56\linewidth}
工場のロボットは、部品をつかんで、決まった場所へ正確に置きます。ロボットの腕には、人の肩やひじのような関節があり、1つ1つの関節をモータで回します。\par\vspace{3pt}
手を「ここ」に置きたいとき、肩とひじを、それぞれ何度回せばよいのでしょう。反対に、関節の角度がわかっているとき、手がどこにあるかは計算できるでしょうか。
\end{minipage}\hfill
\begin{minipage}[c]{0.4\linewidth}\centering
\begin{tikzpicture}[x=0.62mm, y=0.62mm]
  \fill[bkGray!25] (-8,-3.2) rectangle (60,-1.4);
  \pxhbase
  \pxharm{30}{26}{70}{-95}{bkBlue!70}
  \draw[bkInk, line width=0.9pt] ($(P)+(-1.2,-2.6)$) -- ($(P)+(1.2,2.6)$);
  \fill[bkOrange!70] (44,-1.4) rectangle (54,6);
  \draw[bkRed, line width=0.9pt, dashed] (49,10) circle (3.2);
  \draw[bk arrow, bkRed!80!black, line width=0.8pt] (49,24) -- (49,14.4);
  \node[font=\sffamily\footnotesize, text=bkRed!80!black, anchor=south] at (49,24.4) {ここに置きたい};
  \node[font=\sffamily\footnotesize, anchor=east] at ($(E)+(-2.2,0)$) {ひじ ?°};
  \node[font=\sffamily\footnotesize, anchor=west] at (4,4.5) {肩 ?°};
\end{tikzpicture}
\end{minipage}
\end{hatena}

この章では、転がる車輪、ロボットの腕、エンジンのピストン、歯車のように、「回る部品」と「動く部品」が組み合わさった機械の動きを、式で書きます。どれも、平らな面の中で、ずれながら回る動き（\textbf{平面運動}）です。\bkidx{へいめんうんどう}{平面運動}

\begin{yattemiyou}[厚紙の腕で、ねらった点に手を置こう]
\begin{enumerate}[label={\bkcircled[bkGreen]{\arabic*}}, leftmargin=2.4\zw, itemsep=1pt, topsep=2pt]
\item 厚紙を細長く2本切る。穴と穴の間の長さを、\qty{15}{cm}（上の腕）と\qty{12}{cm}（下の腕）にする。
\item 2本を割りピンでつなぐ（ひじ）。上の腕のはしを、方眼紙の点 $(0,\ 0)$ にピンでとめる（肩）。
\item 分度器で、肩の角を\ang{30}、ひじの曲げる角を\ang{45}にして、手の先の点を方眼紙に写そう。
\item 反対に、手の先を点 $(\qty{20}{cm},\ \qty{10}{cm})$ に置いて、肩とひじの角をはかろう。置き方は何通りあるかな。
\end{enumerate}
\end{yattemiyou}

\noindent ③では、手の先は約 $(16.1,\ 19.1)$（cm）の所に来ます。④では、ひじを下にする置き方（肩 約$-3$°、ひじ 約69°）と、ひじを上にする置き方（肩 約57°、ひじ 約$-69$°）の2通りがあります。角から位置を出すのも、位置から角を出すのも、紙とピンで確かめられました。この章では、それを計算でします。

\begin{omoidasou}
\bkrecall{4-7}{2本の腕（2リンク）の手の位置（順運動学）：\\$x=L_1\cos\theta_1+L_2\cos(\theta_1+\theta_2)$\\$y=L_1\sin\theta_1+L_2\sin(\theta_1+\theta_2)$\\手の位置から角を出すときは、三角形OEPの3辺から余弦定理で（逆運動学）。}
\bkrecall{4-2}{弧の長さ $s=r\theta$（$\theta$ はラジアン）。}
\bkrecall{4-8}{半径 $r$ の円を角速度 $\omega$ で回る点の速さ $v=r\omega$。}
\bkrecall{4-14}{平面のモーメント $M=xF_y-yF_x$（左回りが正）。外積 $\vec a\times\vec b$ の向きは、$\vec a$ から $\vec b$ へ回す右ねじの進む向き。}
\bkrecall{力-1}{等速円運動の加速度は中心向きで、大きさ $r\omega^2=\dfrac{v^2}{r}$。}
\bkrecall{力-2}{加速する乗り物の中では、慣性力 $-m\vec A$ を加えると、止まっている人と同じ形の運動方程式で考えられる。}
\bkrecall{力-6}{固定した軸のまわりの回転は $I\dfrac{\dd\omega}{\dd t}=N$。回転のエネルギー $\dfrac{1}{2}I\omega^2$、動力 $P=N\omega$。球 $\dfrac{2}{5}mr^2$、円柱 $\dfrac{1}{2}mr^2$、輪 $mr^2$。}
\bkrecall{7-6}{回転行列 $R(\theta)=\begin{pmatrix}\cos\theta&-\sin\theta\\ \sin\theta&\cos\theta\end{pmatrix}$ をかけると、ベクトルは左回りに $\theta$ 回る。$R(\alpha)R(\beta)=R(\alpha+\beta)$。}
\end{omoidasou}

%% ================================================================
\Needspace{18\baselineskip}
\bksection{doguu}{物理の道具}

\phstage{me}{転がる車輪とロボットの腕を見る}{3}

\bksubsection{1}{転がる車輪の、地面にふれた点は止まっている}
\leavevmode\label{ph:力-7:sec1}自転車で走っているところを、横から写真にとると、ふしぎなことが起きます。車輪の上のほうのスポークはぶれて写り、下のほうのスポークはくっきり写るのです。シャッターが開いているわずかな間に、上のほうは大きく動き、下のほうはほとんど動かないからです。同じ1つの車輪なのに、場所によって速さがちがうのです。

\phboxhead{ペットボトルのキャップを転がしてみる}
キャップのふちに、ペンで点をつけます。机の上で、すべらないように、ゆっくり転がしてみましょう。点は、下の図のような、アーチをつないだ形をえがきます。この曲線を\textbf{サイクロイド}といいます。点が机にふれる所で、曲線はとがっています。点はそこで、一瞬だけ止まっているのです。
\begin{bkfigure}
\begin{tikzpicture}[x=1mm, y=1mm]
  \draw[bkInk, line width=0.9pt] (-4,0) -- (98,0);
  \draw[bkRed, line width=1.3pt, domain=0:12.566, samples=120, variable=\t] plot ({7*(\t-sin(\t r))},{7*(1-cos(\t r))});
  \foreach \t in {0, 3.1416, 6.2832, 9.4248, 12.566} {
    \pxhwheel{({7*\t},7)}{7}{-\t*57.3}
    \fill[bkRed] ({7*(\t-sin(\t r))},{7*(1-cos(\t r))}) circle (1.1);}
  \draw[bk arrow, bkBlue, line width=1pt] (30,20) -- (44,20);
  \node[font=\sffamily\footnotesize, text=bkBlue!80!black, anchor=west] at (45,20) {進む向き};
  \node[font=\sffamily\footnotesize, anchor=north, align=center] at (0,-1.6) {ここで\\止まる};
  \node[font=\sffamily\footnotesize, anchor=north, align=center] at (43.98,-1.6) {ここで\\止まる};
  \node[font=\sffamily\footnotesize, anchor=north, align=center] at (87.96,-1.6) {ここで\\止まる};
\end{tikzpicture}
\bkfigcaption{転がる車輪のふちの点（赤）がえがく線（サイクロイド）。1回転で、円周の長さ $2\pi r$ だけ進む。点が地面にふれた瞬間、点は止まっている。}
\end{bkfigure}

\phboxhead{転がる＝進む＋回る}
転がる車輪の動きを、2つに分けて考えます。車輪全体が右へ速さ $v$ で進む動き（\textbf{並進}）と、車輪が中心のまわりを回る動き（\textbf{回転}）です。回転では、ふちの点は中心のまわりを速さ $r\omega$ で回ります（4-8）。2つの速さを、ベクトルとしてたします（4-10）。
\begin{bkfigure}
\begin{tikzpicture}[x=1mm, y=1mm]
  \foreach \cx/\lbl in {0/並進：全体が右へ $v$, 46/回転：中心のまわりに $r\omega$, 92/合わせると} {
    \draw[bkInk!75, line width=1.3pt] (\cx,0) circle (11);
    \fill[bkGray] (\cx,0) circle (0.9);
    \draw[bkInk, line width=0.8pt] (\cx-15,-11) -- (\cx+15,-11);
    \node[font=\sffamily\footnotesize, anchor=north] at (\cx,-13) {\lbl};}
  % 並進
  \foreach \p in {(0,11),(0,0),(0,-11),(-11,0),(11,0)} {\draw[bk arrow, bkBlue, line width=1.1pt] \p -- ++(8,0);}
  % 回転（右回り：上は右、下は左、右は下、左は上）
  \draw[bk arrow, bkOrange!90!black, line width=1.1pt] (46,11) -- ++(8,0);
  \draw[bk arrow, bkOrange!90!black, line width=1.1pt] (46,-11) -- ++(-8,0);
  \draw[bk arrow, bkOrange!90!black, line width=1.1pt] (57,0) -- ++(0,-8);
  \draw[bk arrow, bkOrange!90!black, line width=1.1pt] (35,0) -- ++(0,8);
  % 合わせると
  \draw[bk arrow, bkRed, line width=1.3pt] (92,11) -- ++(16,0);
  \draw[bk arrow, bkRed, line width=1.3pt] (92,0) -- ++(8,0);
  \draw[bk arrow, bkRed, line width=1.3pt] (103,0) -- ++(8,-8);
  \draw[bk arrow, bkRed, line width=1.3pt] (81,0) -- ++(8,8);
  \fill[bkRed] (92,-11) circle (1.3);
  \node[font=\sffamily\footnotesize, text=bkRed!80!black, anchor=south] at (100,12.2) {$2v$};
  \node[font=\sffamily\footnotesize, text=bkRed!80!black, anchor=east] at (89.6,-8.6) {0};
  \node[font=\sffamily\bfseries\large] at (23,0) {＋};
  \node[font=\sffamily\bfseries\large] at (69,0) {＝};
\end{tikzpicture}
\bkfigcaption{転がる車輪の点の速さ。すべらないとき、回る速さ $r\omega$ は進む速さ $v$ と同じ大きさになる。いちばん上は $2v$、地面にふれた点は0。}
\end{bkfigure}

いちばん上の点では、進む向きと回る向きがそろうので、速さは $v+r\omega$ です。地面にふれた点では向きが反対なので、$v-r\omega$ です。タイヤがすべらないで転がるとき、地面にふれた点は地面に対して動きません。もし動いたら、タイヤは地面をこすって、ブレーキの跡のような線が残るはずです。だから
\[v-r\omega=0\qquad\text{つまり}\qquad v=r\omega\]
です。\label{ph:力-7:korogari}これが\textbf{すべらずに転がる条件}です。いちばん上の点の速さは $v+r\omega=2v$ で、自転車の速さの2倍です。写真で上のスポークがぶれるのは、このためです。

\phboxhead{別の考え方：1回転で進む距離}
タイヤに絵の具をぬって、すべらないように1回転転がすと、地面にはタイヤの円周と同じ長さ $2\pi r$ の跡がつきます。1回転（$2\pi$\,rad）で $2\pi r$ 進むので、1\,rad 回るごとに $r$ 進みます（4-2の $s=r\theta$）。1秒に $\omega$\,rad 回れば、1秒に $r\omega$ 進みます。やはり $v=r\omega$ です。
\pheq{v=r\omega}{m/s ＝ m × rad/s}
\begin{phyomi}[$v=r\omega$]
左の $v$ は、車輪の中心が進む速さ。右は、半径と角速度のかけ算で、ふちの点が中心のまわりを回る速さ。rad は長さ÷長さで単位のない数なので、$\mathrm{m}\times\mathrm{rad/s}=\mathrm{m/s}$ です。この式は、「すべらない」ときだけ成り立ちます。
\end{phyomi}
\begin{phmijika}[自転車の車輪]
半径\qty{0.33}{m}の車輪の自転車が、速さ\qty{5}{m/s}（\qty{18}{km/h}）で走ると、角速度は
\[\omega=5\div0.33\approx\qty{15}{rad/s}\]
です。1回転は $2\pi\approx6.28$\,rad なので、1秒に約2.4回転、毎分約145回転です。いちばん上の点は\qty{10}{m/s}で、車と同じくらいの速さで前へ動いています。
\end{phmijika}
\begin{phbaisuru}{車輪の半径 $r$}
同じ速さで走るなら、$\omega=\dfrac{v}{r}$ なので、回る速さは半分になります。大きな車輪はゆっくり回り、小さな車輪はせわしなく回ります。
\begin{bkcenter}
\begin{tikzpicture}
\begin{axis}[bkplot, width=86mm, height=44mm, xmin=0, xmax=0.75, ymin=0, ymax=55, xtick={0.1,0.2,0.3,0.4,0.5,0.6}, ytick={10,20,30,40,50}, xlabel={$r$（m）}, ylabel={$\omega$（rad/s）}, xticklabel style={/pgf/number format/fixed, /pgf/number format/precision=1}]
  \addplot[bkcurve, domain=0.1:0.7, samples=60] {5/x};
  \fill[bkRed] (axis cs:0.33,15.15) circle (2pt); \fill[bkRed] (axis cs:0.66,7.58) circle (2pt);
  \node[font=\sffamily\footnotesize, anchor=south west] at (axis cs:0.34,16.5) {自転車 $r=0.33$：約15};
  \node[font=\sffamily\footnotesize, anchor=north east] at (axis cs:0.67,6.2) {$r$ が2倍：約7.6};
\end{axis}
\end{tikzpicture}
\bkfigcaption{速さ\qty{5}{m/s}で転がる車輪の、半径と角速度。$r$ と $\omega$ は反比例（2-15）。}
\end{bkcenter}
\end{phbaisuru}

\bksubsection{2}{2本の腕の先の位置は、2つの角で決まる}
\leavevmode\label{ph:力-7:sec2}ロボットの腕のように、硬い棒（\textbf{リンク}）を関節でつないだしくみを、\textbf{リンク機構}といいます。\bkidx{りんくきこう}{リンク機構}ここでは、肩とひじの2つの関節をもつ、平面の上で動く腕（2リンクの腕）を考えます。角の決め方は4-7と同じです。$\theta_1$ は上の腕が $x$ 軸から回った角、$\theta_2$ は下の腕が「上の腕をまっすぐのばした向き」からさらに回った角で、どちらも左回りを正にします。
\begin{bkfigure}
\begin{tikzpicture}[x=0.9mm, y=0.9mm]
  \begin{scope}
  \clip (-46,-6) rectangle (46,46);
  \fill[bkTheme!10, even odd rule] (0,0) circle (44) circle (4);
  \draw[bkTheme!70, line width=0.6pt] (0,0) circle (44) (0,0) circle (4);
  \end{scope}
  \draw[bkGray, line width=0.6pt, -{Stealth[length=4pt]}] (-46,0) -- (48,0) node[anchor=north, font=\sffamily\footnotesize] {$x$};
  \draw[bkGray, line width=0.6pt, -{Stealth[length=4pt]}] (0,-5) -- (0,48) node[anchor=east, font=\sffamily\footnotesize] {$y$};
  \pxharm{24}{20}{58.59}{-71.54}{bkGray!70}
  \pxharm{24}{20}{-5.46}{71.54}{bkBlue!70}
  \node[font=\sffamily\footnotesize, anchor=west] at ($(P)+(1.6,0)$) {P};
  \node[font=\sffamily\footnotesize, text=bkBlue!80!black, anchor=west] at (31,5) {ひじが下};
  \node[font=\sffamily\footnotesize, text=bkGray, anchor=south east] at (12,21.6) {ひじが上};
  \node[font=\sffamily\footnotesize, text=bkDeep, anchor=south] at (-26,38) {手が届く所};
  \node[font=\sffamily\footnotesize, anchor=west, align=left, text width=40mm] at (51,30) {● 2つの角を決めると、\\\ \ 手の位置は1つに決まる。\\[3pt]● 手の位置を決めると、\\\ \ 角の組はふつう2つある。\\[3pt]● 色のついた輪の外には\\\ \ 届かない。};
\end{tikzpicture}
\bkfigcaption{$L_1=\qty{0.30}{m}$、$L_2=\qty{0.25}{m}$ の腕。手を点P $(0.40,\ 0.20)$（m）に置く置き方は、ひじが下（青）とひじが上（灰）の2通り。手が届くのは、肩からの距離が $0.05$〜\qty{0.55}{m} の輪の中（図は上半分）。}
\end{bkfigure}

\begin{phmodel}[2本の腕]
\phignore{腕の太さと重さ。腕のたわみ。関節のすきま（がた）。モータの動きのおくれ。}
\phassume{腕は長さの変わらない硬い棒（剛体）。関節は、平面の中だけで回る。}
\phlimit{重い物を持って腕がたわむとき。速く動かして、モータが角の指令に追いつかないとき。}
\end{phmodel}
図からわかることを、ことばにしておきます。
\begin{enumerate}[label={\bkcircled[bkBlue]{\arabic*}}, leftmargin=2.4\zw, itemsep=1pt, topsep=2pt]
\item 2つの角 $\theta_1$、$\theta_2$ を決めると、手の位置はただ1つに決まります。角から位置を求めることを\textbf{順運動学}といいます（4-7）。
\item 手の位置を決めても、角の組は1つとはかぎりません。ひじを下にする形と上にする形の、ふつう2通りがあります。位置から角を求めることを\textbf{逆運動学}といいます（4-7）。
\item 肩からの距離が、2本をまっすぐのばした $L_1+L_2$ より遠い所や、たたんだ $|L_1-L_2|$ より近い所には、手は届きません。届く所は、輪の形をしています。
\end{enumerate}


%% ================================================================
\phstage{te}{表計算で腕を動かす・歯車を数える}{3}

\bksubsection{3}{表計算で、手をねらった所へ}
\phboxhead{角から位置へ（順運動学）}\label{ph:力-7:sec3}
4-7の式を、表計算に入れます。A列に $\theta_1$、B列に $\theta_2$（度）を入れ、C列とD列に手の $x$、$y$ を計算させます。表計算の \texttt{COS} はラジアンで計算するので、\texttt{RADIANS} で度をラジアンに直します（4-2）。C2には
\[\texttt{=0.30*COS(RADIANS(A2))+0.25*COS(RADIANS(A2+B2))}\]
と入れ、D2にはCOSをSINにかえた式を入れて、下の行へコピーします。
\begin{bkcenter}
\bksheet{4}{26}{%
1 & $\theta_1$（°） & $\theta_2$（°） & $x$（m） & $y$（m） \\
2 & 0 & 45 & 0.477 & 0.177 \\
3 & 15 & 45 & 0.415 & 0.294 \\
4 & 30 & 45 & 0.325 & 0.391 \\
5 & 45 & 45 & 0.212 & 0.462 \\
6 & 60 & 45 & 0.085 & 0.501 \\
7 & 75 & 45 & $-0.047$ & 0.506 \\
8 & 90 & 45 & $-0.177$ & 0.477 \\}
\bktabcaption{$L_1=\qty{0.30}{m}$、$L_2=\qty{0.25}{m}$ の腕の、手の位置。}
\end{bkcenter}
\noindent\begin{minipage}[c]{0.46\linewidth}\centering
\begin{tikzpicture}[x=0.78mm, y=0.78mm]
  \draw[bkGray, line width=0.6pt, -{Stealth[length=4pt]}] (-22,0) -- (52,0) node[anchor=north, font=\sffamily\footnotesize] {$x$};
  \draw[bkGray, line width=0.6pt, -{Stealth[length=4pt]}] (0,-3) -- (0,56) node[anchor=east, font=\sffamily\footnotesize] {$y$};
  \draw[bkOrange, line width=0.7pt, dashed] (20.33:50.88) arc (20.33:110.33:50.88);
  \pxharm{30}{25}{30}{45}{bkBlue!70}
  \foreach \x/\y in {47.7/17.7, 41.5/29.4, 32.5/39.1, 21.2/46.2, 8.5/50.1, -4.7/50.6, -17.7/47.7} {\fill[bkRed] (\x,\y) circle (1.1);}
  \node[font=\sffamily\footnotesize, anchor=north west] at (1.5,-0.8) {O};
\end{tikzpicture}
\end{minipage}\hfill
\begin{minipage}[c]{0.5\linewidth}
表の7つの点（赤）を図にかきました。ひじの角を\ang{45}のまま肩だけを回すと、手は肩を中心とする円（点線）の上を動きます。三角形OEPの形が変わらないので、肩から手までの距離が変わらないからです。\par\vspace{3pt}
ロボットに「まっすぐ動け」と命令するには、肩とひじを同時に、うまく組み合わせて回す必要があります。
\end{minipage}

\phboxhead{位置から角へ（逆運動学）}
\leavevmode\bkhatenaref に答えます。手を置きたい点 $(x,\ y)$ から、4-7の手順で角を求めます。
\begin{enumerate}[label={\bkcircled[bkBlue]{\arabic*}}, leftmargin=2.4\zw, itemsep=1pt, topsep=2pt]
\item 肩から手までの距離 $d=\sqrt{x^2+y^2}$（4-4）。
\item 三角形OEPの3辺 $L_1$、$L_2$、$d$ から、余弦定理でひじの所の角 $\angle\mathrm{OEP}$ を求め、$\theta_2=\ang{180}-\angle\mathrm{OEP}$。
\item OPが $x$ 軸となす角 $\alpha$（$\tan\alpha=\dfrac{y}{x}$）と、OEとOPのなす角 $\beta$（余弦定理）から、$\theta_1=\alpha-\beta$（ひじが下の形）。
\end{enumerate}
この手順をプログラムにして、手を直線 $x=\qty{0.40}{m}$ の上で、$y=0$ から\qty{0.30}{m}まで\qty{0.05}{m}ずつ動かすときの角を計算します。
\begin{phcode}
from math import sqrt, acos, atan2, degrees
L1, L2 = 0.30, 0.25  \# うでの長さ（m）
x = 0.40  \# 手を動かす線 x = 0.40
for i in range(7):
    y = 0.05 * i
    d = sqrt(x**2 + y**2)  \# 肩から手までの距離
    c = (L1**2 + L2**2 - d**2) / (2 * L1 * L2)
    t2 = 180 - degrees(acos(c))  \# ひじの角
    b = (L1**2 + d**2 - L2**2) / (2 * L1 * d)
    t1 = degrees(atan2(y, x)) - degrees(acos(b))  \# 肩の角
    print(round(y, 2), round(t1, 1), round(t2, 1))
\end{phcode}
\begin{phcodeout}
0.0 -38.6 87.1
0.05 -31.1 86.2
0.1 -23.0 83.3
0.15 -14.4 78.5
0.2 -5.5 71.5
0.25 4.1 62.2
0.3 14.5 49.5
\end{phcodeout}
\noindent\begin{minipage}[c]{0.5\linewidth}
\texttt{acos} は cos から角を求める関数（電卓の $\cos^{-1}$）、\texttt{atan2(y, x)} は点 $(x,\ y)$ が $x$ 軸となす角を求める関数です（4-15の極座標の $\theta$）。\texttt{**2} は2乗です。\par\vspace{3pt}
手が上へ動くにつれて、肩の角はふえ、ひじの角はへっていきます。2つの角が、決まった組み合わせで同時に変わると、手はまっすぐ上へ動くのです。
\end{minipage}\hfill
\begin{minipage}[c]{0.46\linewidth}\centering
\begin{tikzpicture}[x=1mm, y=1mm]
  \draw[bkGray, line width=0.6pt, -{Stealth[length=4pt]}] (-4,0) -- (56,0) node[anchor=north, font=\sffamily\footnotesize] {$x$};
  \draw[bkGray, line width=0.6pt, -{Stealth[length=4pt]}] (0,-22) -- (0,36) node[anchor=east, font=\sffamily\footnotesize] {$y$};
  \draw[bkOrange, line width=1pt, dashed] (40,-3) -- (40,33) node[anchor=south, font=\sffamily\footnotesize, text=bkOrange!80!black] {$x=0.40$};
  \foreach \a/\b/\c in {-38.6/87.1/bkBlue!30, -23.0/83.3/bkBlue!45, -5.5/71.5/bkBlue!60, 14.5/49.5/bkBlue!80} {\pxharm{30}{25}{\a}{\b}{\c}}
  \foreach \y in {0,10,20,30} {\fill[bkRed] (40,\y) circle (1.1);}
  \node[font=\sffamily\footnotesize, anchor=west] at (41.5,30) {$y=0.30$};
  \node[font=\sffamily\footnotesize, anchor=north west] at (41,-0.6) {$y=0$};
\end{tikzpicture}
\end{minipage}

\par\vspace{4pt}
\textbf{はてな？の答え}：手を置きたい点から、余弦定理で肩とひじの角を計算します（逆運動学）。答えはふつう2通りあるので、ぶつかる物のないほうを選びます。手を道すじにそって動かすときは、道すじを細かい点に分けて、点ごとにこの計算をくり返し、その角の指令を関節のモータに送ります。工場のロボットは、この計算を、とても短い時間ごとにくり返しています。

\bksubsection{4}{歯車とベルト：歯の数で速さが決まる}
\leavevmode\label{ph:力-7:sec4}モータは、速く回るのは得意ですが、1回転で出せる回す力（トルク）は小さめです。そこで、歯車を使って回る速さを落とし、そのかわりにトルクを大きくします。歯車の働きを、歯を数えて調べます。
\begin{bkfigure}
\begin{tikzpicture}[x=1mm, y=1mm]
  \pxhgear{(0,0)}{8}{20}{0}{bkOrange!55}
  \pxhgear{(32,0)}{24}{60}{183}{bkBlue!35}
  \fill[bkRed] (7.2,0) circle (1.1);
  \fill[bkRed] (8.8,0) circle (1.1);
  \draw[bkInk, line width=0.8pt, ->, >={Stealth[length=4pt]}] (110:12) arc (110:40:12);
  \draw[bkInk, line width=0.8pt, ->, >={Stealth[length=4pt]}] ($(32,0)+(70:28)$) arc (70:110:28);
  \node[font=\sffamily\footnotesize, anchor=north] at (0,-11) {20歯};
  \node[font=\sffamily\footnotesize, anchor=north] at (32,-27) {60歯};
  \node[font=\sffamily\footnotesize, anchor=west, align=left, text width=56mm] at (62,2) {● 小さい歯車が1回転すると、歯が20こ送られる。\\[3pt]● 大きい歯車も20こ分だけ回る。\\\ \ 20÷60＝$\dfrac{1}{3}$ 回転。\\[3pt]● 回る向きは反対。};
\end{tikzpicture}
\bkfigcaption{かみ合う歯車。歯の大きさがそろっているので、歯の数は、かみ合う円の周の長さ（半径）に比例する。}
\label{ph:力-7:fig-gear}
\end{bkfigure}

かみ合っている所では、歯が1つずつ送られます。小さい歯車（20歯）が1回転すると、20この歯が送られ、大きい歯車（60歯）も20こ分、つまり $\dfrac{20}{60}=\dfrac{1}{3}$ 回転します。小さい歯車を毎分3000回転で回すと、大きい歯車は毎分1000回転です。
\[\text{速さの比}\quad\frac{\text{大きい歯車の回る速さ}}{\text{小さい歯車の回る速さ}}=\frac{\text{小さい歯車の歯の数}}{\text{大きい歯車の歯の数}}=\frac{20}{60}\]
回る速さの比は、歯の数の逆の比です。速さを落とす割合 $\dfrac{60}{20}=3$ を\textbf{減速比}といいます。\bkidx{げんそくひ}{減速比}

\phboxhead{トルクはどうなる？}
かみ合う所で、2つの歯車は同じ大きさの力で押し合っています（作用・反作用、力-2）。その力を $F$ とすると、トルク＝力×半径なので、小さい歯車は $F\times r_1$、大きい歯車は $F\times r_2$ です。半径が3倍なので、トルクも3倍になります。モータが\qty{0.5}{N.m}で回せば、大きい歯車の軸からは\qty{1.5}{N.m}が出てきます。
\begin{bkcenter}
\begin{bktable}{colspec={Q[l,m,wd=40mm] Q[c,m,wd=30mm] Q[c,m,wd=22mm] Q[c,m,wd=26mm]}, rowsep=1.5pt}
 & 回る速さ & トルク & 速さ×トルク \\
小さい歯車（モータ） & 毎分3000回転 & \qty{0.5}{N.m} & 1500 \\
大きい歯車 & 毎分1000回転 & \qty{1.5}{N.m} & 1500 \\
\end{bktable}
\bktabcaption{速さは $\dfrac{1}{3}$ になり、トルクは3倍になる。かけ算は変わらない（損失がないとき）。}
\end{bkcenter}
速さとトルクのかけ算が同じになりました。このかけ算は、1秒あたりの仕事、つまり動力に比例します（⑨でくわしく）。歯車は、動力を増やしてはくれません。速さとトルクを、とりかえるだけです。てこや滑車の「仕事の原理」（力-4・力-5）と同じです。

\phboxhead{ベルトとチェーン}
ベルトでつないだ2つの車（プーリ）では、ベルトはすべらずに同じ速さで動きます。ベルトの速さは、どちらのプーリでも「半径×角速度」なので、$r_1\omega_1=r_2\omega_2$。回る速さの比は、半径の逆の比です。自転車のチェーンも同じで、前と後ろの歯車（スプロケット）の歯の数で、ペダルと後ろの車輪の回る速さの比が決まります。ベルトやチェーンでは、2つの車は同じ向きに回ります。
\begin{bkfigure}
\begin{tikzpicture}[x=1mm, y=1mm]
  \draw[bkInk, line width=1pt] ({-6*sin(14.48)},{6*cos(14.48)}) -- ({48-18*sin(14.48)},{18*cos(14.48)});
  \draw[bkInk, line width=1pt] ({-6*sin(14.48)},{-6*cos(14.48)}) -- ({48-18*sin(14.48)},{-18*cos(14.48)});
  \draw[bkInk, line width=1pt] ({-6*sin(14.48)},{6*cos(14.48)}) arc (104.48:255.52:6);
  \draw[bkInk, line width=1pt] ({48-18*sin(14.48)},{18*cos(14.48)}) arc (104.48:-104.48:18);
  \filldraw[fill=bkOrange!40, draw=bkInk, line width=0.6pt] (0,0) circle (5.4);
  \filldraw[fill=bkBlue!25, draw=bkInk, line width=0.6pt] (48,0) circle (17.4);
  \fill[bkInk] (0,0) circle (0.9) (48,0) circle (1.2);
  \draw[bk arrow, bkRed, line width=1pt] (14,9.5) -- (26,12.3);
  \node[font=\sffamily\footnotesize, text=bkRed!80!black, anchor=south] at (20,12.6) {ベルトの速さ $v$};
  \node[font=\sffamily\footnotesize, anchor=north] at (0,-8) {半径 $r_1$};
  \node[font=\sffamily\footnotesize, anchor=north] at (48,-20) {半径 $r_2$};
  \node[font=\sffamily\footnotesize, anchor=west, align=left] at (72,4) {$v=r_1\omega_1=r_2\omega_2$\\[2pt]同じ向きに回る};
\end{tikzpicture}
\bkfigcaption{ベルトでつないだプーリ。ベルトがすべらなければ、どちらのふちも、ベルトと同じ速さ $v$ で動く。}
\end{bkfigure}

\phkokomade

%% ================================================================
\phstage{shiki}{平面運動と回る座標の式から導く}{5}

\bksubsection{5}{剛体の平面運動＝重心の並進＋回転}
\leavevmode\label{ph:力-7:sec5}①では、転がる車輪を「進む＋回る」に分けました。これを、どんな形の硬い物（剛体）にも使える式にします。

\phboxhead{剛体の点の速度}
平面の上を動く剛体の、重心をG、ほかの1つの点をPとします。GからPへのベクトルを $\vec r=\overrightarrow{\mathrm{GP}}$ とします。剛体は形が変わらないので、$\vec r$ の長さ $\rho$ はいつも同じです。変わるのは向きだけで、剛体が $\theta$ 回れば、剛体の上のすべての矢印が同じ $\theta$ だけ回ります。そこで、$\vec r$ の向きの角を $\varphi_0+\theta(t)$ と書くと（$\varphi_0$ ははじめの向き）、
\[\vec r=\rho\bigl(\cos(\varphi_0+\theta),\ \sin(\varphi_0+\theta)\bigr)\]
これを時間で微分します（合成関数の微分、5-6）。$\dfrac{\dd\theta}{\dd t}=\omega$ と書くと、
\[\frac{\dd\vec r}{\dd t}=\rho\omega\bigl(-\sin(\varphi_0+\theta),\ \cos(\varphi_0+\theta)\bigr)=\omega\,(-r_y,\ r_x)\]
これは、$\vec r$ を左回りに\ang{90}回して $\omega$ 倍した矢印です（4-8、4-15の $r\omega\,\vec e_\theta$）。3次元で書くと、$z$ の向きの角速度ベクトル $\vec\omega=(0,\ 0,\ \omega)$ と $\vec r=(r_x,\ r_y,\ 0)$ の外積（4-14）$\vec\omega\times\vec r=(-\omega r_y,\ \omega r_x,\ 0)$ と同じです。Pの位置は「Gの位置＋$\vec r$」なので、微分して
\pheq{\vec v_\mathrm{P}=\vec v_\mathrm{G}+\vec\omega\times\vec r}{m/s ＝ m/s ＋ rad/s × m}
\begin{phyomi}[$\vec v_\mathrm{P}=\vec v_\mathrm{G}+\vec\omega\times\vec r$]
剛体のどの点の速度も、「重心の速度」と「重心のまわりに回る速度」のたし算。回る速度は、GからPへの矢印に直角で、大きさは $\rho\omega$（Gから遠い点ほど速い）。$\vec\omega$ は、回る軸の向きの矢印で、左回りなら紙の表へ向かう（右ねじ、4-14）。
\end{phyomi}
\noindent\begin{minipage}[c]{0.48\linewidth}\centering
\begin{tikzpicture}[x=1mm, y=1mm, rotate=12]
  \filldraw[fill=bkTheme!10, draw=bkTheme!70, line width=0.8pt, rounded corners=5mm] (-14,-9) rectangle (30,13);
  \fill[bkInk] (0,0) circle (1.1);
  \fill[bkInk] (22,6) circle (1.1);
  \draw[bkGray, line width=0.7pt, -{Stealth[length=3pt]}] (0,0) -- (21,5.7);
  \draw[bk arrow, bkBlue, line width=1.3pt] (0,0) -- (13,-3);
  \draw[bk arrow, bkBlue, line width=1pt, dashed] (22,6) -- (35,3);
  \draw[bk arrow, bkOrange!90!black, line width=1.3pt] (22,6) -- (19.6,14.8);
  \draw[bk arrow, bkRed, line width=1.5pt] (22,6) -- (32.6,11.8);
  \node[font=\sffamily\footnotesize, anchor=north east] at (-0.6,-0.6) {G};
  \node[font=\sffamily\footnotesize, anchor=north] at (22,4.6) {P};
  \node[font=\sffamily\footnotesize, text=bkGray, anchor=south] at (9,3.6) {$\vec r$};
  \node[font=\sffamily\footnotesize, text=bkBlue!80!black, anchor=north] at (7,-3.2) {$\vec v_\mathrm{G}$};
  \node[font=\sffamily\footnotesize, text=bkOrange!80!black, anchor=east] at (19.4,13.6) {$\vec\omega\times\vec r$};
  \node[font=\sffamily\footnotesize, text=bkRed!80!black, anchor=west] at (33,12) {$\vec v_\mathrm{P}$};
\end{tikzpicture}
\end{minipage}\hfill
\begin{minipage}[c]{0.48\linewidth}
青の $\vec v_\mathrm{G}$ を点Pに写し（点線）、だいだいの $\vec\omega\times\vec r$ をつなぐと、赤の $\vec v_\mathrm{P}$ になります（4-10の矢印のたし算）。$\vec\omega\times\vec r$ は、$\vec r$ にいつも直角です。
\end{minipage}

\phboxhead{転がる車輪にあてはめる}
半径 $r$ の車輪が右へ転がるとき、車輪は右回り（時計回り）に回るので、$\vec\omega=(0,\ 0,\ -\omega)$ です（$\omega>0$）。重心は中心で、$\vec v_\mathrm{G}=(v,\ 0,\ 0)$。地面にふれた点は $\vec r=(0,\ -r,\ 0)$ の所なので、外積の成分の式（4-14）で
\[\vec\omega\times\vec r=\bigl(0\times0-(-\omega)\times(-r),\ \ (-\omega)\times0-0\times0,\ \ 0\bigr)=(-r\omega,\ 0,\ 0)\]
だから、地面にふれた点の速度は $(v-r\omega,\ 0,\ 0)$。すべらない（この点が止まっている）ので、$v=r\omega$ です。いちばん上の点（$\vec r=(0,\ r,\ 0)$）では $\vec\omega\times\vec r=(r\omega,\ 0,\ 0)$ で、速度は $(v+r\omega,\ 0,\ 0)=(2v,\ 0,\ 0)$。①で図から読んだことが、式から出てきました。

\phboxhead{重心の運動と、重心のまわりの回転}
動きの式は、2つに分けて立てます。剛体の質量を $m$ とします。
\begin{enumerate}[label={\bkcircled[bkBlue]{\arabic*}}, leftmargin=2.4\zw, itemsep=1pt, topsep=2pt]
\item 重心の運動：$m\dfrac{\dd\vec v_\mathrm{G}}{\dd t}=\sum\vec F$。外から加わる力をすべて重心に集めたとして、質点と同じ式で動きます（力-3 の質点系の重心）。
\item 重心のまわりの回転：$I_\mathrm{G}\dfrac{\dd\omega}{\dd t}=\sum N_\mathrm{G}$。$I_\mathrm{G}$ は重心を通る軸のまわりの慣性モーメント、右は重心のまわりのトルクの合計です（力-6 の回転の式）。
\end{enumerate}
\begin{seikaku}
\phboxhead{重心が加速していても、②が使えるわけ}
力-6 の回転の式は、止まっている軸のまわりで作りました。重心は加速しているかもしれないのに、なぜ②が使えるのでしょう。重心といっしょに動く人から見ると、剛体のどの小さな部分にも、その部分の質量に比例した慣性力 $-\Delta m\,\vec A_\mathrm{G}$（力-2）がはたらきます。これは、どの部分にも同じ向きの重力がかかるのと同じ形なので、合計は重心にかかる1つの力と同じです。重心にかかる力は、重心のまわりのトルクを生みません。だから、重心のまわりなら、加速していても、止まっている軸と同じ②が成り立つのです。重心でない点のまわりでは、こうはなりません。
\end{seikaku}

\phboxhead{斜面を転がる物の加速度}
\leavevmode\label{ph:力-7:sec5b}角 $\theta$ の斜面を、半径 $r$、質量 $m$、重心のまわりの慣性モーメント $I$ の丸い物が、すべらずに転がり下ります。
\begin{phmodel}[斜面を転がる物]
\phignore{空気の抵抗。転がり抵抗（物と斜面のわずかな変形でエネルギーが熱になること）。}
\phassume{物も斜面も剛体。すべらずに転がる（$v=r\omega$）。重心は物の中心。}
\phlimit{斜面が急すぎて、すべりが起きるとき。中に液体が入った缶のように、物が剛体でないとき。}
\end{phmodel}
\noindent\begin{minipage}[c]{0.46\linewidth}\centering
\begin{tikzpicture}[x=1mm, y=1mm]
  \fill[bkGray!25] (0,0) -- (52,0) -- (0,30) -- cycle;
  \draw[bkInk, line width=0.9pt] (0,30) -- (52,0) -- (0,0);
  \draw[bkInk, line width=0.5pt] (44,0) arc (180:150:8);
  \node[font=\sffamily\footnotesize] at (39.5,2.6) {$\theta$};
  \coordinate (C) at ($(22,{30-22*30/52})+(60:7.6)$);
  \coordinate (T) at (22,{30-22*30/52});
  \filldraw[fill=bkBlue!25, draw=bkBlue!80!black, line width=0.8pt] (C) circle (7.6);
  \fill[bkInk] (C) circle (0.8);
  \draw[bk arrow, bkRed, line width=1.2pt] (C) -- ++(0,-16);
  \node[font=\sffamily\footnotesize, text=bkRed!80!black, anchor=north] at ($(C)+(0,-16.3)$) {$mg$};
  \draw[bk arrow, bkGreen!70!black, line width=1.2pt] (T) -- ++(60:13);
  \node[font=\sffamily\footnotesize, text=bkGreen!50!black, anchor=south east] at ($(T)+(60:13)+(-0.4,0)$) {$N$};
  \draw[bk arrow, bkOrange!90!black, line width=1.2pt] (T) -- ++(150:10);
  \node[font=\sffamily\footnotesize, text=bkOrange!80!black, anchor=south] at ($(T)+(150:10)+(0,0.4)$) {$f$};
\end{tikzpicture}
\end{minipage}\hfill
\begin{minipage}[c]{0.5\linewidth}
はたらく力は、重力 $mg$、垂直抗力 $N$、斜面にそって上向きの静止摩擦力 $f$ の3つです。斜面にそって下向きを正にし、加速度を $a$、角加速度を $\alpha=\dfrac{\dd\omega}{\dd t}$ とします。\par\vspace{3pt}
重心のまわりのトルクを生むのは、摩擦力だけです（重力は重心に、垂直抗力は重心へ向かってかかるので、うでが0）。
\end{minipage}

\par\vspace{4pt}
①重心の運動（斜面にそって）、②回転、③すべらない条件の3つの式を立てます。
\[\text{①}\ \ ma=mg\sin\theta-f\qquad \text{②}\ \ I\alpha=fr\qquad \text{③}\ \ a=r\alpha\]
③は $v=r\omega$ を時間で微分したものです。②と③から $f=\dfrac{I\alpha}{r}=\dfrac{Ia}{r^2}$。これを①に入れて、
\[ma=mg\sin\theta-\frac{I}{r^2}a\qquad\Bigl(m+\frac{I}{r^2}\Bigr)a=mg\sin\theta\]
両辺を $m$ でわって、\label{ph:力-7:shamen}
\pheq{a=\dfrac{g\sin\theta}{1+\dfrac{I}{mr^2}}}{m/s$^2$ ＝ m/s$^2$}
\begin{phyomi}[$a=\dfrac{g\sin\theta}{1+\frac{I}{mr^2}}$]
分子の $g\sin\theta$ は、すべる物（摩擦なし）の加速度（力-2）。分母の $\dfrac{I}{mr^2}$ は、「回すのにどれだけ手間がかかるか」を表す、単位のない数です。この数の分だけ、加速度が小さくなります。球は $\dfrac{2}{5}$、円柱は $\dfrac{1}{2}$、輪（パイプ）は1です（力-6）。
\end{phyomi}
\begin{phmijika}[ボールとすべり台]
ころころ転がり下りるボールは、氷のように摩擦なしですべり下りる物より、ゆっくり下ります。重力がした仕事の一部が、回転のエネルギーになるからです。
\end{phmijika}
\begin{phbaisuru}{半径 $r$ や質量 $m$}
$I$ は $mr^2$ に比例するので（球なら $\dfrac{2}{5}mr^2$）、$\dfrac{I}{mr^2}$ は変わりません。大きい球も小さい球も、重い球も軽い球も、同じ加速度で転がります。加速度を決めるのは、大きさや重さではなく、「質量が中心からどれだけはなれて分かれているか」という形だけです。
\begin{bkcenter}
\begin{tikzpicture}
\begin{axis}[bkplot, width=92mm, height=46mm, xmin=0, xmax=1.25, ymin=0, ymax=1.15, xtick={0.2,0.4,0.6,0.8,1}, ytick={0.2,0.4,0.6,0.8,1}, xlabel={$\dfrac{I}{mr^2}$}, ylabel={$\dfrac{a}{g\sin\theta}$}, xticklabel style={/pgf/number format/fixed, /pgf/number format/precision=1}, yticklabel style={/pgf/number format/fixed, /pgf/number format/precision=1}]
  \addplot[bkcurve, domain=0:1.2, samples=60] {1/(1+x)};
  \fill[bkRed] (axis cs:0,1) circle (2pt); \fill[bkRed] (axis cs:0.4,0.7143) circle (2pt); \fill[bkRed] (axis cs:0.5,0.6667) circle (2pt); \fill[bkRed] (axis cs:1,0.5) circle (2pt);
  \node[font=\sffamily\footnotesize, anchor=south west] at (axis cs:0.03,1.0) {すべる物（回らない）};
  \node[font=\sffamily\footnotesize, anchor=north east] at (axis cs:0.38,0.69) {球 $\dfrac{5}{7}$};
  \node[font=\sffamily\footnotesize, anchor=south west] at (axis cs:0.52,0.68) {円柱 $\dfrac{2}{3}$};
  \node[font=\sffamily\footnotesize, anchor=south west] at (axis cs:1.0,0.52) {輪 $\dfrac{1}{2}$};
\end{axis}
\end{tikzpicture}
\bkfigcaption{$\dfrac{I}{mr^2}$ が大きい（質量が外側に多い）ほど、ゆっくり転がる。}
\end{bkcenter}
\end{phbaisuru}

\phboxhead{エネルギーで確かめる}
高さ $h$ だけ下りたときの速さを、エネルギー（力-4）で求めてみます。転がる物の運動エネルギーは、重心の動きの $\dfrac{1}{2}mv^2$ と、回転の $\dfrac{1}{2}I\omega^2$ の和です。すべらないので、角速度は速さ $v$ を半径 $r$ でわったものです。これを入れると、
\[mgh=\frac{1}{2}mv^2+\frac{1}{2}I\Bigl(\frac{v}{r}\Bigr)^2=\frac{1}{2}mv^2\Bigl(1+\frac{I}{mr^2}\Bigr)\qquad v^2=\frac{2gh}{1+\frac{I}{mr^2}}\]
斜面の長さを $L$ とすると $h=L\sin\theta$ で、上の加速度の式から作った $v^2=2aL$（力-1）と、同じになります。ここで、摩擦力 $f$ は仕事をしていません。物が斜面にふれている点は、その瞬間に止まっていて（①）、力がかかる点が動かないからです。

\phboxhead{すべらないために必要な摩擦}
①から、摩擦力は $f=mg\sin\theta-ma=mg\sin\theta\cdot\dfrac{I/(mr^2)}{1+I/(mr^2)}$ です。静止摩擦力の最大は $\mu mg\cos\theta$（力-2）なので、すべらないためには、球なら $\mu\geqq\dfrac{2}{7}\tan\theta$ が必要です。\ang{30}の斜面なら $\mu\geqq0.165$ くらいです。これより摩擦が小さいと、球はすべりながら転がります（\bkreidairef{3}）。

\begin{phflow}
\phflowrow{$\vec v_\mathrm{P}=\vec v_\mathrm{G}+\vec\omega\times\vec r$}{地面にふれた点が\\止まっている}{$v=r\omega$}
\phflowrow{$m\dfrac{\dd v}{\dd t}=\sum F$、$I\dfrac{\dd\omega}{\dd t}=\sum N_\mathrm{G}$}{斜面・すべらない\\$a=r\alpha$}{$a=\dfrac{g\sin\theta}{1+\frac{I}{mr^2}}$}
\phflowrow{$a=\dfrac{g\sin\theta}{1+\frac{I}{mr^2}}$}{$I\to0$\\（回らない物）}{$a=g\sin\theta$（力-2）}
\end{phflow}

\phboxhead{地面にふれた点のまわりで考える}
\leavevmode\label{ph:力-7:sec5c}半径 $r$ の丸い物がすべらずに転がるとき、地面にふれた点Qのまわりのトルクで考えると、摩擦力と垂直抗力が式から消えて、便利なことがあります。右向きに転がる向き（時計回り）を正とし、水平な力の合計を $F_x$（摩擦をのぞく）、摩擦力を $f$（右向きが正）とします。重心の運動 $ma=F_x+f$ と、重心のまわりの回転 $I\alpha=N_\mathrm{G}-fr$（$N_\mathrm{G}$ は摩擦をのぞくトルク。摩擦は下のはしで右へ押すので、左回りに回す）と、$a=r\alpha$ から $f$ を消すと、
\[(I+mr^2)\,\alpha=N_\mathrm{G}+F_xr=N_\mathrm{Q}\]
右の $N_\mathrm{Q}$ は、点Qのまわりの（摩擦をのぞく）トルクです。重心はQの真上 $r$ の所なので、水平な力 $F_x$ のトルクが、Gのまわりより $F_xr$ だけ大きくなるからです（4-14）。左の $I+mr^2$ は、点Qのまわりの慣性モーメントです（力-6 の平行軸の定理）。

\begin{motto}[5]
\phboxhead{転がる車輪は、地面にふれた点のまわりを回っている}
転がる車輪のどの点の速度も、地面にふれた点Qから見て「Qからの距離×$\omega$」の大きさで、Qからの線に直角です。だから、その瞬間だけを見ると、車輪全体が点Qのまわりを角速度 $\omega$ で回っていると考えられます。この点を\textbf{瞬間中心}といいます。いちばん上の点はQから $2r$ はなれているので、速さは $2r\omega=2v$。①の図と同じです。機構学では、リンク機構の速さを求めるのに、瞬間中心をよく使います。
\end{motto}

\bksubsection{6}{リンク機構：ピストン・クランクと4節リンク}
\phboxhead{ピストン・クランク}\label{ph:力-7:sec6}
自動車のエンジンでは、燃えたガスがピストンを押し、ピストンの行ったり来たりの動きを、連接棒（コンロッド）とクランクが「回る動き」に変えます。反対に、ポンプやコンプレッサでは、回る動きを行ったり来たりに変えます。
\begin{bkfigure}
\begin{tikzpicture}[x=1mm, y=1mm]
  \draw[bkGray, line width=0.6pt, dashed] (0,0) circle (8);
  \fill[bkGray!25] (24,5.6) rectangle (52,7.2) (24,-5.6) rectangle (52,-7.2);
  \draw[bkInk, line width=0.7pt] (24,5.6) -- (52,5.6) (24,-5.6) -- (52,-5.6);
  \coordinate (A) at (4,6.93);
  \coordinate (B) at (31.13,0);
  \filldraw[fill=bkOrange!45, draw=bkInk, line width=0.7pt] (29.5,-5.4) rectangle (41.5,5.4);
  \draw[bkBlue!80!black, line width=2.4pt, line cap=round] (0,0) -- (A);
  \draw[bkGreen!60!black, line width=2.4pt, line cap=round] (A) -- (B);
  \fill[bkInk] (0,0) circle (1.1);
  \filldraw[fill=white, draw=bkInk] (A) circle (1) (B) circle (1);
  \draw[bkInk, line width=0.5pt] (5,0) arc (0:60:5);
  \node[font=\sffamily\footnotesize] at (30:7.6) {$\theta$};
  \node[font=\sffamily\footnotesize, text=bkBlue!80!black, anchor=east] at (1.4,4.8) {$r$};
  \node[font=\sffamily\footnotesize, text=bkGreen!50!black, anchor=south] at (18.5,4.2) {$l$};
  \draw[bkInk, line width=0.6pt, {Stealth[length=3pt]}-{Stealth[length=3pt]}] (0,-12) -- (31.13,-12);
  \node[font=\sffamily\footnotesize, anchor=north] at (15.6,-12.6) {$x$};
  \draw[bkGray, line width=0.4pt] (31.13,-1.2) -- (31.13,-13) (0,-1.2) -- (0,-13);
  \node[font=\sffamily\footnotesize, anchor=north] at (0,-15.6) {O（クランクの軸）};
  \node[font=\sffamily\footnotesize, anchor=south] at (35.5,7.6) {ピストン};
  \node[font=\sffamily\footnotesize, anchor=west, align=left] at (58,5) {Aからおろした垂線で\\直角三角形を作ると\\$x=r\cos\theta+\sqrt{l^2-r^2\sin^2\theta}$};
\end{tikzpicture}
\bkfigcaption{ピストン・クランク。クランクの半径 $r$、連接棒の長さ $l$。クランクが角 $\theta$ のときのピストンの位置 $x$。}
\end{bkfigure}
クランクのはしAの高さは $r\sin\theta$、横の位置は $r\cos\theta$ です。Aから $x$ 軸へ垂線をおろすと、連接棒 $l$ を斜辺とする直角三角形ができます。三平方の定理（4-4）で、横の辺は $\sqrt{l^2-r^2\sin^2\theta}$。だから
\pheq{x=r\cos\theta+\sqrt{l^2-r^2\sin^2\theta}}{m ＝ m}
\begin{phyomi}[ピストンの位置]
1つめの項は、クランクのはしの横の位置。2つめの項は、連接棒が横に占める長さで、連接棒がかたむくほど短くなる。$\theta=\ang{0}$ でいちばん右（$l+r$）、$\theta=\ang{180}$ でいちばん左（$l-r$）。ピストンが動くはば $2r$ を\textbf{行程}（ストローク）といいます。
\end{phyomi}
\begin{bkcenter}
\begin{tikzpicture}
\begin{axis}[bkplot, width=112mm, height=48mm, xmin=0, xmax=380, ymin=-48, ymax=52, xtick={90,180,270,360}, ytick={-40,-20,20,40}, xlabel={$\theta$（°）}, ylabel={$x-l$（mm）}]
  \addplot[bkcurve, domain=0:360, samples=121] {40*cos(x)+sqrt(140^2-40^2*sin(x)^2)-140};
  \addplot[bkcurve2, dashed, domain=0:360, samples=121] {40*cos(x)};
  \node[font=\sffamily\footnotesize, text=bkBlue!80!black, anchor=west] at (axis cs:150,38) {青：ピストン};
  \node[font=\sffamily\footnotesize, text=bkOrange!80!black, anchor=west] at (axis cs:150,24) {点線：$r\cos\theta$};
\end{axis}
\end{tikzpicture}
\bkfigcaption{$r=\qty{40}{mm}$、$l=\qty{140}{mm}$ のときのピストンの位置（青）。点線の $r\cos\theta$ より、\ang{90}のとき\qty{5.8}{mm}だけ左（行程のまん中より下死点の側）にいる。}
\end{bkcenter}
ピストンは、$r\cos\theta$ の波（4-8）に近い動きをしますが、少しくずれています。連接棒が長いほど（$l$ が $r$ に比べて大きいほど）、ずれは小さくなります。

ピストンの速さは、$x$ を時間で微分して求めます。$\theta=\omega t$ とし、合成関数の微分（5-6）を使うと、
\[\frac{\dd x}{\dd t}=-r\omega\sin\theta\Bigl(1+\frac{r\cos\theta}{\sqrt{l^2-r^2\sin^2\theta}}\Bigr)\]
です。1回転の間にピストンは行程を2回（行きと帰り）動くので、毎秒 $n$ 回転なら、\textbf{平均の速さ}は $2\times2r\times n=4rn$ です。
\begin{phbaisuru}{回転数 $n$}
平均の速さも、いちばん速いときの速さも、2倍になります。毎分3000回転・$r=\qty{40}{mm}$ なら、平均は\qty{8}{m/s}、いちばん速いのは $\theta\approx\ang{75}$ のときで約\qty{13}{m/s}です。ピストンは1秒に100回も、止まっては動きを変えています。
\end{phbaisuru}

\phboxhead{4節リンク}
4本の棒をつないで輪にしたしくみを\textbf{4節リンク}といいます。1本（固定リンク）を動かないようにとめ、もう1本（クランク）を回すと、向かいの1本（てこ）が、決まったはばで行ったり来たりします。自動車のワイパのしくみに使われてきました。
\begin{bkfigure}
\begin{tikzpicture}[x=1mm, y=1mm]
  \fill[bkTheme!12] (36,0) -- ++(93.58:24) arc (93.58:153.62:24) -- cycle;
  \draw[bkGray, line width=0.6pt, dashed] (0,0) circle (12);
  \fill[bkGray!50] (-4,-4.4) rectangle (40,-2);
  \draw[bkInk, line width=1.6pt] (0,0) -- (36,0);
  \coordinate (A) at (6,10.39);
  \coordinate (B) at (32.83,23.79);
  \draw[bkBlue!80!black, line width=2.4pt, line cap=round] (0,0) -- (A);
  \draw[bkGreen!60!black, line width=2.4pt, line cap=round] (A) -- (B);
  \draw[bkOrange!90!black, line width=2.4pt, line cap=round] (36,0) -- (B);
  \filldraw[fill=white, draw=bkInk] (0,0) circle (1.2) (36,0) circle (1.2) (A) circle (1) (B) circle (1);
  \node[font=\sffamily\footnotesize, anchor=north] at (0,-5) {$\mathrm{A_0}$};
  \node[font=\sffamily\footnotesize, anchor=north] at (36,-5) {$\mathrm{B_0}$};
  \node[font=\sffamily\footnotesize, anchor=south east] at (A) {A};
  \node[font=\sffamily\footnotesize, anchor=south west] at (B) {B};
  \node[font=\sffamily\footnotesize, text=bkBlue!80!black, anchor=west] at (2.6,3.4) {$a$};
  \node[font=\sffamily\footnotesize, text=bkGreen!50!black, anchor=south east] at (19,17.4) {$b$};
  \node[font=\sffamily\footnotesize, text=bkOrange!80!black, anchor=west] at (35.2,12) {$c$};
  \node[font=\sffamily\footnotesize, anchor=south] at (18,0.4) {$d$};
  \draw[bkInk, line width=0.5pt] (42,0) arc (0:97.6:6);
  \node[font=\sffamily\footnotesize] at ($(36,0)+(40:8.4)$) {$\psi$};
  \draw[bkInk, line width=0.5pt] (4,0) arc (0:60:4);
  \node[font=\sffamily\footnotesize] at (30:6.4) {$\varphi$};
  \node[font=\sffamily\footnotesize, anchor=west, align=left, text width=52mm] at (54,14) {クランク $a=\qty{2}{cm}$\\つなぎ $b=\qty{5}{cm}$\\てこ $c=\qty{4}{cm}$\\固定 $d=\qty{6}{cm}$\\[3pt]クランクが1回転すると、てこは約\ang{94}〜約\ang{154}の間（色の部分、約\ang{60}）を往復する。};
\end{tikzpicture}
\bkfigcaption{4節リンク（クランク・てこ機構）。てこの角 $\psi$ は、三角形 $\mathrm{AB_0B}$ で余弦定理を使って求める（\bkreidairef{10}）。}
\end{bkfigure}
クランクがぐるっと1回転できるかどうかは、4本の長さで決まります。いちばん短い棒の長さ $s$ と、いちばん長い棒の長さ $l$ の和が、ほかの2本の和以下（$s+l\leqq p+q$）なら、いちばん短い棒は1回転できます（\textbf{グラスホフの条件}）。上の図では $2+6=8\leqq5+4=9$ で、いちばん短いクランクが回れます。

\bksubsection{7}{腕を行列で書く：回転行列・同次変換・ヤコビ行列}
\phboxhead{回転行列で書く}\label{ph:力-7:sec7}
上の腕は、$x$ 軸の向きの矢印 $\begin{pmatrix}L_1\\0\end{pmatrix}$ を $\theta_1$ 回したものです。下の腕は、$\begin{pmatrix}L_2\\0\end{pmatrix}$ を $\theta_2$ 回し、さらに上の腕といっしょに $\theta_1$ 回したものです。回転行列（7-6）で書くと、
\[\vec p=R(\theta_1)\begin{pmatrix}L_1\\0\end{pmatrix}+R(\theta_1)R(\theta_2)\begin{pmatrix}L_2\\0\end{pmatrix}\]
$R(\theta_1)R(\theta_2)=R(\theta_1+\theta_2)$ なので（7-6）、成分にすると4-7の順運動学の式になります。関節が3つ、4つとふえても、行列を1つずつかけていけばよいのです。

\phboxhead{同次変換：回して、進む}
「回して、進む」を1つの行列にまとめる方法があります（7-7）。2次元の点 $(x,\ y)$ に1を加えた $(x,\ y,\ 1)$ を使い、
\[A(\theta,\ L)=\begin{pmatrix}\cos\theta&-\sin\theta&L\cos\theta\\ \sin\theta&\cos\theta&L\sin\theta\\ 0&0&1\end{pmatrix}\]
と決めます。左上の2×2が回転、右の列が「回った向きに $L$ 進む」移動です。この行列を\textbf{同次変換行列}といいます。肩からひじまでが $A(\theta_1,\ L_1)$、ひじから手までが $A(\theta_2,\ L_2)$ なので、肩から手までは、かけ算 $A(\theta_1,\ L_1)A(\theta_2,\ L_2)$ の1つの行列になります。その右の列の上の2つが、手の位置 $(x,\ y)$ です。\bkreidairef{7}で計算してみます。6つの関節をもつ産業用ロボットでは、3次元の4×4の同次変換行列を6つかけて、手の位置と向きを求めます。

\phboxhead{逆運動学の式}
4-7の余弦定理の手順を、式にまとめておきます。$\cos\angle\mathrm{OEP}=-\cos\theta_2$ なので、
\[\cos\theta_2=\frac{x^2+y^2-L_1^2-L_2^2}{2L_1L_2}\]
この式を満たす $\theta_2$ は、＋と－の2つあります（ひじが下と上）。右辺が1より大きいか $-1$ より小さいときは、手が届かず、答えはありません。$\theta_1$ は、OPの向きの角 $\alpha$ から、OEとOPのなす角 $\beta$ を引く（または足す）ことで求めます。$\beta$ は、$\tan\beta=\dfrac{L_2\sin\theta_2}{L_1+L_2\cos\theta_2}$ からも求められます。

\phboxhead{手の速さ：ヤコビ行列}
関節を回す速さ $\dot\theta_1=\dfrac{\dd\theta_1}{\dd t}$、$\dot\theta_2=\dfrac{\dd\theta_2}{\dd t}$ がわかっているとき、手の速さはいくらでしょう。順運動学の式を時間で微分します（合成関数の微分、5-6）。$\theta_{12}=\theta_1+\theta_2$ と書くと、
\[\dot x=-L_1\sin\theta_1\,\dot\theta_1-L_2\sin\theta_{12}\,(\dot\theta_1+\dot\theta_2)\]
\[\dot y=L_1\cos\theta_1\,\dot\theta_1+L_2\cos\theta_{12}\,(\dot\theta_1+\dot\theta_2)\]
$\dot\theta_1$ と $\dot\theta_2$ でまとめて、行列で書くと（7-2）、
\pheq{\begin{pmatrix}\dot x\\ \dot y\end{pmatrix}=\begin{pmatrix}-L_1\sin\theta_1-L_2\sin\theta_{12}&-L_2\sin\theta_{12}\\ L_1\cos\theta_1+L_2\cos\theta_{12}&L_2\cos\theta_{12}\end{pmatrix}\begin{pmatrix}\dot\theta_1\\ \dot\theta_2\end{pmatrix}}{m/s ＝ m × rad/s}
この行列を\textbf{ヤコビ行列} $J$ といいます（偏微分を並べた行列。8-3 で学ぶ）。1列めは「肩だけを回したときの手の速度」、2列めは「ひじだけを回したときの手の速度」です。
\begin{phyomi}[$\dot{\vec p}=J\dot{\vec\theta}$]
関節の角速度に行列 $J$ をかけると、手の速度が出る。$J$ の中身は、そのときの角で変わる。1列めは、肩から手への矢印 $(x,\ y)$ を\ang{90}回した $(-y,\ x)$ になっている。肩を回すと、手は肩を中心とする円の接線の向きに動くからです（4-15）。
\end{phyomi}
\begin{phbaisuru}[0]{ひじの角 $\theta_2$}
$J$ の行列式（7-4）を計算すると、$\det J=L_1L_2\sin\theta_2$ になります。腕がまっすぐのびる（$\theta_2\to0$）と、行列式は0に近づきます。このとき、2つの列が同じ向きになり、手は「腕にそった向き」には、どう関節を回しても動かせません。手を少し動かすだけのために、関節をとても速く回さなければならなくなります。このような姿勢を\textbf{特異姿勢}といい、ロボットを動かすときは、近づかないようにします。
\begin{bkcenter}
\begin{tikzpicture}[x=1mm, y=1mm]
  \draw[bkGray, line width=0.6pt] (-3,0) -- (62,0);
  \pxharm{30}{25}{20}{0}{bkBlue!70}
  \draw[bk arrow, bkRed, line width=1.2pt] (P) -- ++(110:9);
  \draw[bkGray, line width=1pt, dashed] (P) -- ++(20:9);
  \node[font=\sffamily\footnotesize, text=bkRed!80!black, anchor=south] at ($(P)+(110:9.5)$) {動かせる};
  \node[font=\sffamily\footnotesize, text=bkGray, anchor=west] at ($(P)+(20:9.6)$) {動かせない};
\end{tikzpicture}
\bkfigcaption{腕がまっすぐのとき（$\theta_2=\ang{0}$）。肩を回しても、ひじを回しても、手は赤い向きにしか動かない。}
\end{bkcenter}
\end{phbaisuru}

\bksubsection{8}{回る部屋から見る：遠心力とコリオリの力}
\leavevmode\label{ph:力-7:sec8}回転台（メリーゴーラウンド）に乗ると、外へ放り出されそうな力を感じます。台の上でボールを転がすと、ボールはまっすぐ進まず、曲がっていくように見えます。回る物の上から見ると、なぜこんなことが起きるのでしょう。
\begin{bkfigure}
\begin{tikzpicture}[x=1mm, y=1mm]
  \foreach \cx in {0,66} {\filldraw[fill=bkTheme!8, draw=bkTheme!70, line width=0.8pt] (\cx,0) circle (26); \fill[bkInk] (\cx,0) circle (0.9);}
  \draw[bkInk, line width=0.7pt, ->, >={Stealth[length=4pt]}] (20:29) arc (20:60:29);
  \node[font=\sffamily\footnotesize, anchor=south west] at (45:29.6) {$\omega$};
  \draw[bkRed, line width=1.3pt] (0,0) -- (26,0);
  \foreach \s in {0.25,0.5,0.75,1} {\fill[bkRed] ({26*\s},0) circle (1);}
  \draw[bkRed, line width=1.3pt, domain=0:1, samples=40, variable=\s] plot ({66+26*\s*cos(-90*\s)},{26*\s*sin(-90*\s)});
  \foreach \s in {0.25,0.5,0.75,1} {\fill[bkRed] ({66+26*\s*cos(-90*\s)},{26*\s*sin(-90*\s)}) circle (1);}
  \node[font=\sffamily\footnotesize, anchor=north] at (0,-28) {地面から見る：まっすぐ};
  \node[font=\sffamily\footnotesize, anchor=north] at (66,-28) {台から見る：右へ曲がる};
  \node[font=\sffamily\footnotesize, anchor=south, align=center] at (13,1.2) {ボール};
\end{tikzpicture}
\bkfigcaption{左回りの回転台の中心から、摩擦のないボールを外へ転がす。ふちに着くまでに台が\ang{90}回るとき、同じ時間ごとのボールの位置（赤い点）。}
\end{bkfigure}
地面から見ると、ボールには横向きの力がはたらかないので、まっすぐ進みます（慣性の法則、力-2）。でも、その間に台が左へ回るので、台の上の人から見ると、ボールは右へ曲がっていきます。台の上の人がこの動きを「力」で説明しようとすると、ボールを横へ押す力があることになります。これが\textbf{コリオリの力}です。\bkidx{こりおりのちから}{コリオリの力}

\phboxhead{回る座標の運動方程式}
台といっしょに角速度 $\omega$ で回る座標で、物の位置を $\vec r\,'=(x',\ y')$ とします。地面の座標での位置 $\vec r$ は、$\vec r\,'$ を角 $\theta=\omega t$ だけ回したものなので、回転行列（7-6）で $\vec r=R(\theta)\,\vec r\,'$ です。$\omega$ は一定とします。

$R(\theta)$ の4つの成分を $\theta$ で微分すると、
\[\frac{\dd R}{\dd\theta}=\begin{pmatrix}-\sin\theta&-\cos\theta\\ \cos\theta&-\sin\theta\end{pmatrix}=R(\theta)\,R_{90},\qquad R_{90}=\begin{pmatrix}0&-1\\1&0\end{pmatrix}\]
$R_{90}$ は、\ang{90}回す回転行列 $R(\ang{90})$ です（4-8の $(x,\ y)\to(-y,\ x)$）。\ang{90}を2回回すと\ang{180}なので、$R_{90}R_{90}=-E$（$E$ は単位行列）です。積の微分（5-6）で $\vec r=R\,\vec r\,'$ を2回微分すると、$\dfrac{\dd R}{\dd t}=\omega R R_{90}$ なので、
\[\frac{\dd\vec r}{\dd t}=R\Bigl(\frac{\dd\vec r\,'}{\dd t}+\omega R_{90}\vec r\,'\Bigr),\qquad \frac{\dd^2\vec r}{\dd t^2}=R\Bigl(\frac{\dd^2\vec r\,'}{\dd t^2}+2\omega R_{90}\frac{\dd\vec r\,'}{\dd t}-\omega^2\vec r\,'\Bigr)\]
地面の座標では、運動方程式 $m\dfrac{\dd^2\vec r}{\dd t^2}=\vec F$ が成り立ちます（力-2）。両辺に左から $R(-\theta)$ をかけて（回す前にもどす）、回る座標の力 $\vec F\,'=R(-\theta)\vec F$ で書くと、
\pheq{m\frac{\dd^2\vec r\,'}{\dd t^2}=\vec F\,'+m\omega^2\vec r\,'-2m\omega R_{90}\vec v\,'}{N ＝ N ＋ kg × (rad/s)$^2$ × m}
右の2つめと3つめの項が、回る座標の人だけが感じる\textbf{見かけの力}（慣性力）です。
\begin{phyomi}[回る座標の運動方程式]
2つめの $m\omega^2\vec r\,'$ は、中心から外向きで、大きさ $mr\omega^2$（$r$ は中心からの距離）。これが\textbf{遠心力}です。3つめは、速度 $\vec v\,'$ を\ang{90}回して逆向きにした向き、つまり、左回りの台なら進む向きの右へ向かう力で、大きさ $2m\omega v'$。これがコリオリの力です。止まっている物（$\vec v\,'=\vec 0$）には、コリオリの力ははたらきません。3次元のベクトルで書くと、遠心力は $-m\vec\omega\times(\vec\omega\times\vec r\,')$、コリオリの力は $-2m\vec\omega\times\vec v\,'$ です。
\end{phyomi}
回る座標の中で止まって見える物（$\vec v\,'=\vec 0$、$\dfrac{\dd^2\vec r\,'}{\dd t^2}=\vec 0$）では、$\vec F\,'+m\omega^2\vec r\,'=\vec 0$ です。\label{ph:力-7:enshin}
\pheq{\text{遠心力の大きさ}\quad mr\omega^2=\frac{mv^2}{r}}{N ＝ kg × m × (rad/s)$^2$}
地面から見ると、同じ物は半径 $r$ の等速円運動をしていて、力 $\vec F$ は中心向きの $mr\omega^2$（向心力、力-1・力-2）です。台の上から見ると、その力 $\vec F$ と、外向きの遠心力がつり合って止まっている、と説明します。同じ動きを、2通りの見方で説明しているのです。
\begin{phmijika}[洗濯機の脱水]
力-1 の例題の洗濯機（毎分800回転、半径\qty{0.25}{m}）では、$\omega=800\times2\pi\div60\approx\qty{83.8}{rad/s}$、$r\omega^2\approx\qty{1750}{m/s^2}$ で、約\qty{1800}{m/s^2}（$g$ の約180倍）です。いっしょに回る服の上から見ると、\qty{1}{g}の水のつぶに約\qty{1.8}{N}の遠心力（\qty{180}{g}の物の重さくらい）がかかり、水は服からはなれて、かべの穴から外へ飛んでいきます。
\end{phmijika}
\begin{phbaisuru}{角速度 $\omega$}
遠心力は $\omega^2$ に比例するので、4倍になります。回転数を2倍にすると、遠心力は4倍です。同じ半径を回る物の速さ $v$ で言いかえても、$\dfrac{mv^2}{r}$ なので、速さ2倍で4倍です。カーブを速く曲がると急に危なくなるのは、このためです。
\begin{bkcenter}
\begin{tikzpicture}
\begin{axis}[bkplot, width=86mm, height=42mm, xmin=0, xmax=1100, ymin=0, ymax=3800, xtick={200,400,600,800,1000}, ytick={1000,2000,3000}, xlabel={毎分の回転数}, ylabel={$r\omega^2$（m/s$^2$）}]
  \addplot[bkcurve, domain=0:1000, samples=60] {0.25*(x*2*pi/60)^2};
  \fill[bkRed] (axis cs:400,438.6) circle (2pt); \fill[bkRed] (axis cs:800,1754.6) circle (2pt);
  \node[font=\sffamily\footnotesize, anchor=north west] at (axis cs:420,380) {毎分400回転：約440};
  \node[font=\sffamily\footnotesize, anchor=south east] at (axis cs:770,1800) {毎分800回転：約1750};
\end{axis}
\end{tikzpicture}
\bkfigcaption{半径\qty{0.25}{m}の脱水槽のかべでの $r\omega^2$。回転数が2倍で4倍。}
\end{bkcenter}
\end{phbaisuru}

\phboxhead{地球も回っている}
地球は西から東へ、1恒星日（\qty{86164}{s}）に1回転しています。角速度は $\omega=\dfrac{2\pi}{\num{86164}}\approx\qty{7.29e-5}{rad/s}$ です。地面の上の私たちは、回る部屋の中にいるのと同じです。北緯 $\phi$ の地面は、真上を向いた軸のまわりに、$\omega\sin\phi$ の速さで回っています（北極なら $\omega$、赤道なら0）。だから、地面にそって速さ $v$ で動く物には、水平な向きに、大きさ $2m\omega v\sin\phi$ のコリオリの力がはたらきます。北半球では、進む向きの右へ曲げる向きです（\bkreidairef{15}）。
\begin{bkfigure}
\begin{tikzpicture}[x=1mm, y=1mm]
  \fill[bkBlue!10] (0,0) circle (24);
  \node[font=\sffamily\bfseries\large, text=bkBlue!80!black] at (0,0) {低};
  \foreach \a in {0,90,180,270} {
    \draw[bkGray, line width=0.7pt, dashed, -{Stealth[length=3pt]}] (\a:24) -- (\a:7);
    \draw[bkRed, line width=1.3pt, -{Stealth[length=4pt]}] (\a:24) .. controls (\a:15) and (\a+35:11) .. (\a+75:9);}
  \node[font=\sffamily\footnotesize, anchor=west, align=left, text width=62mm] at (32,4) {点線：気圧の低い中心へ吸いこまれる向き。\\赤：北半球では、コリオリの力で、進む向きの右へ少しずつ曲がる。\\全体として、左回りのうずになる（台風・低気圧）。};
\end{tikzpicture}
\bkfigcaption{北半球の台風が左回りのうずになるわけ（上から見た図）。南半球では右回りになる。}
\end{bkfigure}

\bksubsection{9}{歯車列とベルト：速度比とトルク比}
\leavevmode\label{ph:力-7:sec9}④で歯を数えて調べたことを、式で導きます。

\phboxhead{速度比}
2つの歯車がかみ合う点では、歯が押し合いながら動くので、すべりがないとすると、2つの歯車のかみ合う点の速さは同じです。半径を $r_1$、$r_2$、角速度を $\omega_1$、$\omega_2$ とすると、
\[r_1\omega_1=r_2\omega_2\]
歯車をかみ合わせるには、歯の大きさ（ピッチ）がそろっていなければなりません。円周 $2\pi r$ に歯が $z$ こ、同じ間かくで並ぶので、$2\pi r=z\times(\text{ピッチ})$ で、半径は歯の数 $z$ に比例します。だから
\pheq{\frac{\omega_2}{\omega_1}=\frac{r_1}{r_2}=\frac{z_1}{z_2}}{（単位のない数）＝（単位のない数）}

\phboxhead{トルク比と動力}
回る物の動力は $P=N\omega$ です（力-6）。歯車の中で失われる動力がなければ、入る動力と出る動力は等しく、$N_1\omega_1=N_2\omega_2$。だから
\pheq{\frac{N_2}{N_1}=\frac{\omega_1}{\omega_2}=\frac{z_2}{z_1}}{（単位のない数）＝（単位のない数）}
④では、かみ合う点で押し合う力（作用・反作用）から同じ答えを出しました。2通りの考え方で、同じ結果になります。実際の歯車では、摩擦などで少し動力が失われます。出る動力÷入る動力を\textbf{効率} $\eta$ といい（力-4）、出るトルクは $\eta\times\dfrac{z_2}{z_1}N_1$ です。よくできた歯車の1組の効率は、0.95より大きいくらいです。
\begin{phflow}
\phflowrow{かみ合う点の速度が等しい\\$r_1\omega_1=r_2\omega_2$}{歯の大きさがそろう\\$r\propto z$}{$\dfrac{\omega_2}{\omega_1}=\dfrac{z_1}{z_2}$}
\phflowrow{動力 $P=N\omega$}{損失がない\\$N_1\omega_1=N_2\omega_2$}{$\dfrac{N_2}{N_1}=\dfrac{z_2}{z_1}$}
\phflowrow{ベルトの速さ $v=r_1\omega_1=r_2\omega_2$}{ベルトが\\すべらない}{$\dfrac{\omega_2}{\omega_1}=\dfrac{r_1}{r_2}$}
\end{phflow}
\noindent 歯車をいくつもつないだもの（\textbf{歯車列}）では、それぞれの組の速度比をかけ算します。20歯→60歯で $\dfrac{1}{3}$、同じ軸についた15歯→45歯でさらに $\dfrac{1}{3}$ なら、全体で $\dfrac{1}{9}$ です（\bkreidairef{11}）。かみ合うたびに回る向きが反対になるので、2組なら、出口はモータと同じ向きに回ります。

\bksubsection{10}{こまとジャイロの入口：歳差運動}
\leavevmode\label{ph:力-7:sec10}速く回っているこまは、かたむいても倒れず、軸がゆっくりと円をえがいて回ります。自転車の車輪の軸の片方を、ひもでつるしても、車輪が速く回っていれば、落ちずに、水平のままゆっくり回ります。このような軸の回り方を\textbf{歳差運動}といいます。\bkidx{さいさうんどう}{歳差運動}
\begin{bkfigure}
\begin{tikzpicture}[x=1mm, y=1mm]
  \draw[bkInk, line width=0.8pt] (-6,36) -- (6,36);
  \draw[bkInk, line width=0.7pt] (0,36) -- (0,2);
  \draw[bkInk, line width=1.6pt] (0,2) -- (26,2);
  \filldraw[fill=bkBlue!15, draw=bkBlue!80!black, line width=1.2pt] (22,2) ellipse (4 and 14);
  \draw[bkBlue!80!black, line width=0.6pt] (22,2) ellipse (1.6 and 5.6);
  \draw[bk arrow, bkGreen!70!black, line width=1.6pt] (26,2) -- (44,2);
  \node[font=\sffamily\footnotesize, text=bkGreen!50!black, anchor=south] at (40,2.8) {$\vec L=I\omega$};
  \draw[bk arrow, bkRed, line width=1.3pt] (22,-12.6) -- (22,-22);
  \node[font=\sffamily\footnotesize, text=bkRed!80!black, anchor=west] at (23,-19) {$mg$};
  \draw[bk arrow, bkOrange!90!black, line width=1.4pt] (44,2) -- (50,10);
  \node[font=\sffamily\footnotesize, text=bkOrange!80!black, anchor=south, align=center] at (50,10.5) {$\Delta\vec L=\vec N\Delta t$\\（紙の奥へ）};
  \draw[bkInk, line width=0.7pt, ->, >={Stealth[length=4pt]}] (-14,26) arc (200:340:14 and 4);
  \node[font=\sffamily\footnotesize, anchor=east] at (-14,28) {$\Omega$};
  \node[font=\sffamily\footnotesize, anchor=north] at (11,1.2) {$d$};
  \node[font=\sffamily\footnotesize, anchor=west, align=left, text width=50mm] at (64,-6) {● 重力のトルク $\vec N$ は、軸に直角で水平（紙の奥の向き）。\\[2pt]● $\vec L$ は、$\vec N$ の向きに少しずつ向きを変え、軸がひものまわりを回る。};
\end{tikzpicture}
\bkfigcaption{ひもでつるした、速く回る車輪。ひもから車輪の中心までの長さ $d$。重力 $mg$ が、ひもの点のまわりにトルク $N=mgd$ を生む。}
\end{bkfigure}
車輪の角運動量は、軸の向きで大きさ $L=I\omega$ です（力-6）。重力 $mg$ は、ひもの点のまわりに、大きさ $N=mgd$ のトルクを生みます。その向きは、うで（軸の向き）にも重力（下向き）にも直角なので、水平で、軸に直角な向きです（外積、4-14）。

角運動量の式 $\dfrac{\dd\vec L}{\dd t}=\vec N$（力-6）から、短い時間 $\Delta t$ に $\vec L$ は $\vec N\Delta t$ だけ変わります。$\vec N$ は $\vec L$ に直角なので、$\vec L$ の長さは変わらず、向きだけが水平に回ります。角 $\Delta\phi=\dfrac{N\Delta t}{L}$ だけ回るので、軸が回る角速度は
\pheq{\Omega=\frac{N}{L}=\frac{mgd}{I\omega}}{rad/s ＝ N·m ÷ (kg·m$^2$ × rad/s)}
\begin{phyomi}[$\Omega=\dfrac{mgd}{I\omega}$]
重力のトルクが大きいほど、軸は速く回る。車輪の回る速さ $\omega$ が大きいほど、軸はゆっくり回る。ただし、この式は、車輪が十分速く回り、$\Omega$ が $\omega$ よりずっと小さいときの近似です。
\end{phyomi}
\begin{phbaisuru}{車輪の角速度 $\omega$}
$\Omega$ は $\omega$ に反比例するので、半分になります。$m=\qty{2.0}{kg}$、$d=\qty{0.20}{m}$、$I=\qty{0.10}{kg.m^2}$ の車輪では、$\omega=\qty{30}{rad/s}$ で $\Omega\approx\qty{1.3}{rad/s}$、$\omega=\qty{60}{rad/s}$ で約\qty{0.65}{rad/s}。車輪が速く回るほど、軸の向きは変わりにくくなります。これが、ジャイロ（回転を使って向きを保つ道具）の働きです。
\end{phbaisuru}

\begin{phsaki}{3次元の回転とロボットの動力学}
こまの軸がふらふらとゆれる動き（章動）や、3次元で自由に回る物の動きは、オイラーの運動方程式で調べます（慣性モーメントは3×3の行列になる。7-8）。また、ロボットの腕をねらった速さで動かすのに、関節のモータがどれだけのトルクを出せばよいかは、腕全体の運動方程式（動力学）で求めます。関節がいくつもある腕では、ラグランジュの運動方程式（波-3）を使うのが便利です。どちらも、大学の機械力学・ロボット工学で学びます。
\end{phsaki}

\begin{phmada}[自転車は、なぜ倒れずに走れる？]
止まった自転車はすぐ倒れますが、走っている自転車は、手をはなしても倒れにくくなります。この理由は19世紀から研究され、車輪のジャイロの効果や、前の車輪がハンドルの軸より後ろで地面にふれていること（キャスタの効果）などが考えられてきました。2011年、オランダとアメリカの研究者たちは、ジャイロの効果もキャスタの効果もない特別な自転車を作り、それでも、ある速さの範囲では自分で倒れずに立ち直ることを示しました。自転車が安定する理由は、1つの効果ではなく、重さの配置やハンドルの形などがかかわり合って決まることがわかってきました。どんな形の自転車なら乗りやすいのかは、人の運転とあわせて、今も研究が続いています。
\end{phmada}
\begin{phmada}[やわらかい物をつかむロボット]
形の決まった硬い部品なら、この章の計算でロボットの手を正確に動かせます。でも、布・ひも・食品のように、つかむと形が変わる物を、人のように器用にあつかうことは、まだうまくできていません。物がどう変形するかを予想しながら、指の力をどう調節すればよいかは、今も研究が続いています。
\end{phmada}

\begin{senmonsho}[機構学・機械力学の書き方]
\begin{bkhonbun}
剛体上の任意の点Pの速度は $\symbf v_\mathrm{P}=\symbf v_\mathrm{G}+\symbf\omega\times\symbf r_{\mathrm{P/G}}$ で与えられる。角速度 $\symbf\omega$ で回転する座標系から見た質点の運動方程式は $m\symbf a'=\symbf F-2m\symbf\omega\times\symbf v'-m\symbf\omega\times(\symbf\omega\times\symbf r')$ であり、右辺第2項をコリオリ力、第3項を遠心力という。
\end{bkhonbun}
\bkyaku{太字はベクトル（4-13）。$\symbf r_{\mathrm{P/G}}$ は「Gから見たPの位置」で、この章の $\vec r=\overrightarrow{\mathrm{GP}}$。ダッシュ（$'$）は、回る座標で見た量。「右辺第2項」は右の2つめの式のこと。}
\end{senmonsho}

\begin{taisetsu}
\def\tsl#1{\par\noindent\hangindent=6.2\zw\hangafter=1\makebox[6.2\zw][l]{\sffamily\bfseries\small\color{bkDeep}#1}}
\tsl{剛体の点}$\vec v_\mathrm{P}=\vec v_\mathrm{G}+\vec\omega\times\vec r$
\tsl{転がる}$v=r\omega$（上のはしは $2v$、地面にふれた点は0）
\tsl{斜面}$a=\dfrac{g\sin\theta}{1+I/(mr^2)}$
\tsl{2本の腕}順運動学は4-7の式。逆運動学は余弦定理で、答えは2通り。
\tsl{歯車}$\dfrac{\omega_2}{\omega_1}=\dfrac{z_1}{z_2}$、$\dfrac{N_2}{N_1}=\dfrac{z_2}{z_1}$（動力 $N\omega$ は同じ）
\tsl{遠心力}$mr\omega^2$（外向き）
\tsl{コリオリ}$2m\omega v'$（左回りの座標なら、進む向きの右）
\tsl{歳差運動}$\Omega=\dfrac{mgd}{I\omega}$
\end{taisetsu}

\bkkigocards{\bkkigocard{$\vec\omega$}{オメガ（ベクトル）}{角速度ベクトル。大きさは角速度、向きは回る軸の向き（右ねじ）}{angular velocity vector}{力-7}}{\bkkigocard{$J$}{ジェイ}{ヤコビ行列。関節の角速度から手の速度を出す行列。モータの本では慣性モーメントにも使う}{Jacobian matrix}{力-7}}
\begin{kotoba}\small
\bkword{平面運動}{平面の中で、ずれながら回る剛体の動き。重心の並進＋重心のまわりの回転。}
\bkword{転がりの条件}{すべらずに転がるとき、地面にふれた点は止まっている。$v=r\omega$。}
\bkword{順運動学・逆運動学}{関節の角から手の位置を求めること・手の位置から関節の角を求めること。}
\bkword{リンク機構}{硬い棒（リンク）を関節でつないだしくみ。ピストン・クランク、4節リンクなど。}
\bkword{減速比}{歯車などで、回る速さを落とす割合。トルクは同じ割合で大きくなる。}
\bkword{遠心力・コリオリの力}{回る座標の中で見るときだけに加える見かけの力（慣性力）。}
\bkword{歳差運動}{回る物の軸が、トルクを受けて、ゆっくり向きを変えていく動き。}
\end{kotoba}

%% ================================================================
\bksection{reidai}{例題でたしかめよう}

\begin{phreidai}[自転車の車輪の速さ]\label{ph:力-7:ex1}
半径\qty{0.33}{m}の車輪の自転車が、\qty{5.0}{m/s}で走っています。車輪の角速度と、車輪のいちばん上の点・中心・地面にふれた点の、地面に対する速さを求めましょう。
\phbamen
自転車は、まっすぐな道を一定の速さで走っています。車輪は道の上を転がっています。
\phmushi
タイヤのへこみと、タイヤと道のすべりを無視します。車輪は半径\qty{0.33}{m}の円とします。
\phzu
車輪の中心の速さを $v$、角速度を $\omega$、半径を $r$ とします。①の「並進＋回転」の図で考えます。
\phriyu
すべらずに転がるので、地面にふれた点の速さは0で、$v=r\omega$ が使えます。上の点は、進む速さと回る速さの向きがそろいます。
\phtateru
\[\omega=\frac{v}{r}=\frac{5.0}{0.33},\qquad \text{上の点}=v+r\omega\]
\phtoku
\[\omega=15.15\cdots\approx\qty{15}{rad/s},\qquad \text{上の点}=5.0+5.0=\qty{10}{m/s}\]
中心は\qty{5.0}{m/s}、地面にふれた点は0です。
\phtashikame
単位：m/s÷m＝1/s で、rad/s になっています。\\
大きさ：$15\div2\pi\approx2.4$ で、1秒に約2.4回転。自転車をこぐ感じに合います。\\
極端な場合：止まっている（$v=0$）なら、$\omega=0$ で、どの点も止まっています。
\kotae 約\qty{15}{rad/s}。上の点\qty{10}{m/s}、中心\qty{5.0}{m/s}、地面にふれた点0
\end{phreidai}
\begin{jibun}
半径\qty{0.30}{m}のタイヤの車が、\qty{72}{km/h}（\qty{20}{m/s}）で走っています。タイヤの角速度と、タイヤのいちばん上の点の速さを求めましょう。また、タイヤは毎分約何回転していますか。
\bkjans{$\omega=\dfrac{20}{0.30}\approx\qty{67}{rad/s}$。上の点は\qty{40}{m/s}。毎分 $67\times60\div2\pi$ で約640回転}
\end{jibun}

\begin{phreidai}[斜面の転がり競争]\label{ph:力-7:ex2}
角\ang{30}、長さ\qty{2.0}{m}の斜面の上から、球・円柱・輪（うすいパイプ）を同時にそっと放します。それぞれが下に着くまでの時間と、下での速さを求めましょう。どれが先に着きますか。
\phbamen
3つとも、止まった状態から、すべらずに転がり下ります。
\phmushi
空気の抵抗と、転がり抵抗を無視します。3つとも中身のつまり方は一様（輪はうすい）とします。
\phzu
\noindent\begin{minipage}[c]{0.42\linewidth}\centering
\begin{tikzpicture}[x=1mm, y=1mm]
  \fill[bkGray!25] (0,0) -- (40,0) -- (0,23.1) -- cycle;
  \draw[bkInk, line width=0.8pt] (0,23.1) -- (40,0) -- (0,0);
  \draw[bkInk, line width=0.5pt] (33,0) arc (180:150:7);
  \node[font=\sffamily\footnotesize] at (29.4,2.2) {\ang{30}};
  \foreach \x/\c in {4/bkBlue!50, 10/bkOrange!60, 16/bkGreen!50} {
    \filldraw[fill=\c, draw=bkInk, line width=0.5pt] ($(\x,{23.1-\x*23.1/40})+(60:2.6)$) circle (2.6);}
  \node[font=\sffamily\footnotesize, anchor=south west] at (9,22) {球・円柱・輪};
\end{tikzpicture}
\end{minipage}\hfill
\begin{minipage}[c]{0.55\linewidth}
斜面にそって下向きを正、加速度 $a$。長さ $L=\qty{2.0}{m}$、高さ $h=L\sin\ang{30}=\qty{1.0}{m}$。$k=\dfrac{I}{mr^2}$ は、球0.4、円柱0.5、輪1。
\end{minipage}
\phriyu
すべらずに転がるので、⑤の $a=\dfrac{g\sin\theta}{1+k}$ が使えます。加速度が一定なので、力-1 の $L=\dfrac{1}{2}at^2$、$v=at$ が使えます。
\phtateru
\[a=\frac{9.8\times\sin\ang{30}}{1+k}=\frac{4.9}{1+k},\qquad t=\sqrt{\frac{2L}{a}},\qquad v=at\]
\phtoku
\begin{bkcenter}
\begin{bktable}{colspec={Q[l,m,wd=12mm] Q[c,m,wd=10mm] Q[c,m,wd=30mm] Q[c,m,wd=17mm] Q[c,m,wd=18mm]}, rowsep=1.2pt, colsep=3pt}
 & $k$ & $a$（m/s$^2$） & $t$（s） & $v$（m/s） \\
球 & 0.4 & $4.9\div1.4=3.5$ & 1.07 & 3.7 \\
円柱 & 0.5 & $4.9\div1.5=3.27$ & 1.11 & 3.6 \\
輪 & 1 & $4.9\div2=2.45$ & 1.28 & 3.1 \\
\end{bktable}
\end{bkcenter}
\phtashikame
単位：$\sqrt{\text{m}\div\text{m/s}^2}=\text{s}$。\\
別の方法：エネルギーで $v=\sqrt{\dfrac{2gh}{1+k}}$ からも求められます。球なら
\[v=\sqrt{\frac{2\times9.8\times1.0}{1.4}}=\sqrt{14}\approx\qty{3.7}{m/s}\]
で、同じです。\\
極端な場合：$k\to0$（回らない、摩擦なしですべる物）なら $a=\qty{4.9}{m/s^2}$、$t\approx\qty{0.90}{s}$ で、どれよりも速い。
\kotae 球 約\qty{1.1}{s}・\qty{3.7}{m/s}、円柱 約\qty{1.1}{s}・\qty{3.6}{m/s}、輪 約\qty{1.3}{s}・\qty{3.1}{m/s}。球が先に着く
\end{phreidai}
\begin{jibun}
角\ang{20}、長さ\qty{3.0}{m}の斜面で、球と輪を同時にそっと放します。それぞれの加速度と、下に着くまでの時間を、小数第2位まで求めましょう（$\sin\ang{20}=0.342$）。
\bkjans{球：$a=\dfrac{9.8\times0.342}{1.4}\approx\qty{2.39}{m/s^2}$、$t=\sqrt{\dfrac{6.0}{2.39}}\approx\qty{1.58}{s}$。輪：$a\approx\qty{1.68}{m/s^2}$、$t\approx\qty{1.89}{s}$}
\end{jibun}

\begin{phreidai}[ボウリングの球：すべりから転がりへ]\label{ph:力-7:ex3}
ボウリングの球（一様な球とする）を、回さずに\qty{8.0}{m/s}で床に投げこみました。はじめ球は床をすべり、動摩擦力で速さがへり、回る速さがふえて、やがてすべらずに転がるようになります。動摩擦係数を0.20として、（1）転がり始めるまでの時間、（2）そのときの速さ、（3）それまでに進む距離、（4）失われる運動エネルギーの割合を求めましょう。
\phbamen
はじめ $\omega=0$ なので、床にふれた点は\qty{8.0}{m/s}で床をこすります。こする向きと反対の、後ろ向きの動摩擦力がはたらきます。
\phmushi
空気の抵抗、床の油のむら、球の指穴を無視します。動摩擦係数は一定の0.20とします（実際のレーンは油がぬってあり、もっと小さい所が多い）。
\phzu
\noindent\begin{minipage}[c]{0.38\linewidth}\centering
\begin{tikzpicture}[x=1mm, y=1mm]
  \draw[bkInk, line width=0.9pt] (-14,0) -- (20,0);
  \filldraw[fill=bkBlue!20, draw=bkBlue!80!black, line width=0.8pt] (0,8) circle (8);
  \fill[bkInk] (0,8) circle (0.8);
  \draw[bk arrow, bkBlue, line width=1.2pt] (0,8) -- (12,8);
  \node[font=\sffamily\footnotesize, text=bkBlue!80!black, anchor=south] at (8,8.4) {$v$};
  \draw[bk arrow, bkOrange!90!black, line width=1.2pt] (0,0) -- (-9,0);
  \node[font=\sffamily\footnotesize, text=bkOrange!80!black, anchor=north] at (-6,-0.6) {$\mu'mg$};
  \draw[bkInk, line width=0.7pt, ->, >={Stealth[length=3pt]}] (100:11) ++(0,8) arc (100:40:11);
  \node[font=\sffamily\footnotesize, anchor=south] at (8,17.4) {$\omega$};
\end{tikzpicture}
\end{minipage}\hfill
\begin{minipage}[c]{0.6\linewidth}
進む向きを正、球の速さ $v$、右回りの角速度 $\omega$、質量 $m$、半径 $r$、$I=\dfrac{2}{5}mr^2$。動摩擦力は $\mu'mg$（$\mu'=0.20$）で後ろ向き。
\end{minipage}
\phriyu
重心の運動と重心のまわりの回転を、別々に立てます（⑤）。摩擦力は重心を通らないので、球を右回りに回すトルク $\mu'mg\times r$ を生みます。転がり始めるのは、$v=r\omega$ になったときです。
\phtateru
\[m\frac{\dd v}{\dd t}=-\mu'mg,\qquad I\frac{\dd\omega}{\dd t}=\mu'mgr\]
はじめ $v=8.0$、$\omega=0$ なので、時間 $t$ のあと（力-1 の等加速度と同じ形）、
\[v=8.0-\mu'gt,\qquad r\omega=\frac{\mu'mgr^2}{I}t=\frac{5}{2}\mu'gt\]
\phtoku
（1）$v=r\omega$ とおくと、$8.0-\mu'gt=\dfrac{5}{2}\mu'gt$ より $\dfrac{7}{2}\mu'gt=8.0$。
\[t=\frac{2\times8.0}{7\times0.20\times9.8}=\frac{16}{13.72}\approx\qty{1.17}{s}\]
（2）$v=8.0-0.20\times9.8\times1.166\approx\qty{5.71}{m/s}$（$\dfrac{5}{7}\times8.0$ と同じ）。\\
（3）$s=8.0t-\dfrac{1}{2}\mu'gt^2=8.0\times1.166-0.98\times1.166^2\approx\qty{8.0}{m}$。\\
（4）はじめの運動エネルギーは $\dfrac{1}{2}m\times8.0^2=32m$。転がり始めた後は、
\[\frac{1}{2}mv^2+\frac{1}{2}I\omega^2=\frac{1}{2}mv^2\times\frac{7}{5}=0.7m\times5.71^2\approx22.9m\]
失われた割合は
\[\frac{32m-22.9m}{32m}\approx0.29\qquad\Bigl(\text{ちょうど}\ \frac{2}{7}\Bigr)\]
\phtashikame
別の方法：床にふれた点のまわりでは、摩擦力のうでが0なので、この点のまわりの角運動量は変わりません。
\[mv_0r=mvr+I\omega=mvr\times\frac{7}{5}\qquad v=\frac{5}{7}v_0\]
同じになりました。\\
極端な場合：摩擦が大きいほど早く転がり始めますが、転がり始めた速さ（$v_0$ の $\tfrac{5}{7}$ 倍）は、摩擦係数によりません。\\
転がり始めた後は、ふれた点が止まるので摩擦はなくなり（転がり抵抗を無視すれば）、速さは一定です。
\kotae （1）約\qty{1.2}{s}（2）約\qty{5.7}{m/s}（3）約\qty{8.0}{m}（4）$\dfrac{2}{7}$（約29\,\%）
\end{phreidai}
\begin{jibun}
一様な球を、回さずに\qty{6.0}{m/s}で床に投げこみました。動摩擦係数を0.25として、転がり始めるまでの時間、そのときの速さ、それまでに進む距離を求めましょう。
\bkjans{$t=\dfrac{2\times6.0}{7\times0.25\times9.8}\approx\qty{0.70}{s}$。$v=\dfrac{5}{7}\times6.0\approx\qty{4.3}{m/s}$。距離は $s=v_0t-\dfrac{1}{2}\mu'gt^2$ に入れて約\qty{3.6}{m}}
\end{jibun}

\begin{phreidai}[糸巻きを引くと、どちらへ転がる？]\label{ph:力-7:ex4}
質量\qty{0.10}{kg}、外側の半径 $R=\qty{0.050}{m}$、糸を巻く軸の半径 $r=\qty{0.020}{m}$、中心のまわりの慣性モーメント $I=\qty{1.0e-4}{kg.m^2}$ の糸巻きが、床の上にあります。軸の下側から出た糸を、（1）水平に\qty{0.30}{N}で引くとき、どちらへ、どんな加速度で転がりますか。（2）床から\ang{80}の向きに引くと、どちらへ転がりますか。
\phbamen
糸巻きは床の上で、すべらずに転がり始めます。糸を引く向きによって、転がる向きが変わるかもしれません。
\phmushi
糸の質量とのび、空気の抵抗を無視します。床の摩擦は十分大きく、すべらないとします。
\phzu
\noindent\begin{minipage}[c]{0.4\linewidth}\centering
\begin{tikzpicture}[x=0.8mm, y=0.8mm]
  \draw[bkInk, line width=0.9pt] (-16,0) -- (34,0);
  \filldraw[fill=bkOrange!25, draw=bkInk, line width=0.8pt] (0,12) circle (12);
  \filldraw[fill=bkOrange!55, draw=bkInk, line width=0.6pt] (0,12) circle (4.8);
  \fill[bkInk] (0,12) circle (0.8);
  \draw[bkRed, line width=1.1pt] (0,7.2) -- (30,7.2);
  \draw[bk arrow, bkRed, line width=1.2pt] (22,7.2) -- (30,7.2);
  \node[font=\sffamily\footnotesize, text=bkRed!80!black, anchor=south] at (26,7.6) {$F$};
  \fill[bkInk] (0,0) circle (1);
  \node[font=\sffamily\footnotesize, anchor=north] at (0,-0.8) {Q};
  \draw[bkGray, line width=0.6pt, {Stealth[length=3pt]}-{Stealth[length=3pt]}] (-14,0) -- (-14,7.2);
  \node[font=\sffamily\footnotesize, anchor=east, align=right] at (-14.6,3.6) {$R-r$};
\end{tikzpicture}
\end{minipage}\hfill
\begin{minipage}[c]{0.58\linewidth}
床にふれた点をQとします。右向きに転がる向き（時計回り）を正とし、角加速度 $\alpha$、加速度 $a=R\alpha$。水平な糸は、Qから高さ $R-r$ の所を右へ引きます。
\end{minipage}
\phriyu
床の摩擦力の大きさはわからないので、摩擦力が式に出てこない「Qのまわりのトルク」の式 $(I+mR^2)\alpha=N_\mathrm{Q}$（⑤）を使います。重力と垂直抗力は、Qを通る線の上にあるので、Qのまわりのトルクは0です。
\phtateru
（1）$N_\mathrm{Q}=F\times(R-r)$（時計回り）。
\[(I+mR^2)\,\alpha=F(R-r)\]
\phtoku
\[I+mR^2=1.0\times10^{-4}+0.10\times0.050^2=3.5\times10^{-4}\ \text{kg·m}^2\]
\[\alpha=\frac{0.30\times0.030}{3.5\times10^{-4}}\approx\qty{25.7}{rad/s^2}\]
\[a=R\alpha=0.050\times25.7\approx\qty{1.3}{m/s^2}\]
正なので、糸を引く向き（右）へ転がります。\\
（2）糸がななめになると、糸の延長の線とQとの位置関係が変わります。糸の向きの角を $\phi$ とします。糸は軸の円にふれているので、中心Gから糸の線までの距離は $r$ です。GからQまで（長さ $R$）を、糸に直角な向きにはかると $R\cos\phi$ です。だから、Qから糸の線までの距離（うで）は $R\cos\phi-r$ で、Qのまわりのトルクは $F(R\cos\phi-r)$（時計回りが正）です。
\[\phi=\ang{80}:\quad R\cos\phi=0.050\times0.174=0.0087<r=0.020\]
トルクは負なので、左へ（糸を引く向きと反対へ）転がり、糸は巻きとられます。
\phtashikame
単位：N·m÷kg·m$^2$＝1/s$^2$（rad/s$^2$）。\\
向きの確かめ：糸の延長の線がQの上を通れば右、下を通れば左。ちょうどQを通るのは $\cos\phi=\dfrac{r}{R}=0.4$、$\phi\approx\ang{66}$ のときで、そのとき糸巻きは転がり始めません。\\
（1）の摩擦力：$ma=F+f$ から $f=0.10\times1.29-0.30\approx\qty{-0.17}{N}$（左向き）。$\mu\geqq0.17$ くらいあれば、すべりません。
\kotae （1）引く向き（右）へ、約\qty{1.3}{m/s^2}（2）引く向きと反対（左）へ転がる
\end{phreidai}
\begin{jibun}
質量\qty{0.20}{kg}、外側の半径\qty{0.040}{m}、軸の半径\qty{0.015}{m}、中心のまわりの慣性モーメント\qty{1.6e-4}{kg.m^2}の糸巻きの、軸の下側から出た糸を、（1）水平に\qty{0.50}{N}で引くときの加速度を求めましょう。（2）床から\ang{75}の向きに引くと、どちらへ転がりますか。
\bkjans{（1）$I+mR^2=4.8\times10^{-4}$、$\alpha=\dfrac{0.50\times0.025}{4.8\times10^{-4}}\approx26.0$、$a=0.040\times26.0\approx\qty{1.0}{m/s^2}$（引く向き）。（2）$0.040\cos\ang{75}\approx0.010<0.015$ なので、引く向きと反対へ転がる}
\end{jibun}

\begin{phreidai}[2本の腕：角から手の位置]\label{ph:力-7:ex5}
上の腕 $L_1=\qty{0.30}{m}$、下の腕 $L_2=\qty{0.25}{m}$ のロボットの腕で、$\theta_1=\ang{30}$、$\theta_2=\ang{45}$ のとき、ひじと手の位置を求めましょう。
\phbamen
平面の上で動く2リンクの腕です。肩を原点にします。
\phmushi
腕の太さ・たわみ・関節のすきまを無視します（②のモデル化の札）。
\phzu
角の決め方は②と同じ（$\theta_2$ は上の腕の向きからの角）。ひじE、手P。
\phriyu
角から位置を求める順運動学なので、4-7の式を使います。下の腕の向きは $\theta_1+\theta_2=\ang{75}$ です。
\phtateru
\[\mathrm{E}=(0.30\cos\ang{30},\ 0.30\sin\ang{30})\]
\[\mathrm{P}=\mathrm{E}+(0.25\cos\ang{75},\ 0.25\sin\ang{75})\]
\phtoku
$\cos\ang{30}=0.866$、$\sin\ang{30}=0.5$、$\cos\ang{75}=0.259$、$\sin\ang{75}=0.966$ なので、
\[\mathrm{E}=(0.260,\ 0.150)\]
\[\mathrm{P}=(0.260+0.065,\ 0.150+0.241)=(0.325,\ 0.391)\]
\phtashikame
単位：m。\\
大きさ：肩から手までの距離 $\sqrt{0.325^2+0.391^2}\approx0.509$ は、$L_1+L_2=0.55$ より小さく、届く範囲です。③の表の4行めとも同じです。\\
極端な場合：$\theta_2=0$（まっすぐ）なら、手は $(0.55\cos\ang{30},\ 0.55\sin\ang{30})$。
\kotae ひじ $(0.26,\ 0.15)$\,m、手 約 $(0.32,\ 0.39)$\,m
\end{phreidai}
\begin{jibun}
$L_1=\qty{0.40}{m}$、$L_2=\qty{0.30}{m}$ の腕で、$\theta_1=\ang{20}$、$\theta_2=\ang{50}$ のときの手の位置を、小数第3位まで求めましょう。
\bkjans{ひじ $(0.376,\ 0.137)$。手は、ひじから\ang{70}の向きに\qty{0.30}{m}進んで、$(0.478,\ 0.419)$\,m}
\end{jibun}

\begin{phreidai}[2本の腕：手の位置から角（答えは2つ）]\label{ph:力-7:ex6}
$L_1=\qty{0.30}{m}$、$L_2=\qty{0.25}{m}$ の腕で、手を点 $(0.40,\ 0.20)$（m）に置くための $\theta_1$、$\theta_2$ を、2通り求めましょう。
\phbamen
手を置きたい点が決まっていて、関節の角を求める逆運動学の問題です。
\phmushi
腕の太さ・たわみ、関節の角の範囲の制限を無視します。
\phzu
肩O、ひじE、手Pの三角形OEPで考えます（②の図）。$d=\mathrm{OP}$、OPが $x$ 軸となす角 $\alpha$、OEとOPのなす角 $\beta$。
\phriyu
三角形OEPの3辺がわかるので、余弦定理で角が求められます（4-7、③の手順）。
\phtateru
\[d=\sqrt{0.40^2+0.20^2}\]
\[\cos\angle\mathrm{OEP}=\frac{L_1^2+L_2^2-d^2}{2L_1L_2},\qquad \cos\beta=\frac{L_1^2+d^2-L_2^2}{2L_1d}\]
\phtoku
$d=\sqrt{0.20}\approx0.447$。
\[\cos\angle\mathrm{OEP}=\frac{0.09+0.0625-0.20}{0.15}\approx-0.317,\qquad \angle\mathrm{OEP}\approx\ang{108.5}\]
\[\theta_2=\ang{180}-\ang{108.5}=\ang{71.5}\]
$\tan\alpha=\dfrac{0.20}{0.40}=0.5$ より $\alpha\approx\ang{26.6}$。
\[\cos\beta=\frac{0.09+0.20-0.0625}{2\times0.30\times0.447}\approx0.848,\qquad \beta\approx\ang{32.0}\]
ひじが下：$\theta_1=\alpha-\beta\approx\ang{-5.5}$、$\theta_2\approx\ang{71.5}$。\\
ひじが上：$\theta_1=\alpha+\beta\approx\ang{58.6}$、$\theta_2\approx\ang{-71.5}$。
\phtashikame
順運動学で確かめる：ひじが下なら、$\theta_1+\theta_2=\ang{66.0}$ なので、
\[x=0.30\cos(\ang{-5.5})+0.25\cos\ang{66.0}\approx0.299+0.102=0.401\]
\[y=0.30\sin(\ang{-5.5})+0.25\sin\ang{66.0}\approx-0.029+0.228=0.199\]
ほぼ $(0.40,\ 0.20)$ です。\\
届くか：$d=0.447$ は $0.05$〜$0.55$ の間なので、届きます。
\kotae $(\theta_1,\ \theta_2)\approx(\ang{-5.5},\ \ang{71.5})$ と $(\ang{58.6},\ \ang{-71.5})$
\end{phreidai}
\begin{jibun}
$L_1=\qty{0.40}{m}$、$L_2=\qty{0.30}{m}$ の腕で、手を点 $(0.50,\ 0.30)$（m）に置くための角を、2通り求めましょう。
\bkjans{$d\approx0.583$、$\cos\angle\mathrm{OEP}=-0.375$、$\theta_2=\pm\ang{68.0}$。$\alpha\approx\ang{31.0}$、$\beta\approx\ang{28.5}$。$(\ang{2.5},\ \ang{68.0})$ と $(\ang{59.5},\ \ang{-68.0})$}
\end{jibun}

\begin{phreidai}[同次変換行列のかけ算で手の位置]\label{ph:力-7:ex7}
\bkreidairef{5}の腕（$L_1=0.30$、$L_2=0.25$、$\theta_1=\ang{30}$、$\theta_2=\ang{45}$）の手の位置を、同次変換行列 $A(\theta_1,\ L_1)A(\theta_2,\ L_2)$ の右の列から求めましょう。
\phbamen
関節が多いロボットで使う、行列をかけていく方法を、2関節で練習します。
\phmushi
\bkreidairef{5}と同じです。
\phzu
⑦の $A(\theta,\ L)$。数は小数第3位までで計算します。
\phriyu
同次変換行列の積の右の列の上の2つが、手の位置です（⑦）。
\phtateru
\[A(\ang{30},\ 0.30)=\begin{pmatrix}0.866&-0.5&0.260\\0.5&0.866&0.150\\0&0&1\end{pmatrix}\]
\[A(\ang{45},\ 0.25)=\begin{pmatrix}0.707&-0.707&0.177\\0.707&0.707&0.177\\0&0&1\end{pmatrix}\]
\phtoku
積の右の列だけを計算します（7-3：1つめの行列の行と、2つめの行列の右の列の内積）。
\[x=0.866\times0.177+(-0.5)\times0.177+0.260\times1\]
\[\phantom{x}\approx0.065+0.260=0.325\]
\[y=0.5\times0.177+0.866\times0.177+0.150\times1\]
\[\phantom{y}\approx0.242+0.150=0.392\]
\phtashikame
\bkreidairef{5}の $(0.325,\ 0.391)$ とほぼ同じです（小さなちがいは、途中で丸めたため）。\\
積の左上の2×2は $R(\ang{75})$ になり、手の向きが\ang{75}であることもわかります。
\kotae 約 $(0.32,\ 0.39)$\,m
\end{phreidai}
\begin{jibun}
$L_1=\qty{0.40}{m}$、$L_2=\qty{0.30}{m}$、$\theta_1=\ang{90}$、$\theta_2=\ang{-90}$ のとき、$A(\theta_1,\ L_1)A(\theta_2,\ L_2)$ を計算して、手の位置を求めましょう。図もかいて確かめましょう。
\bkjans{$A(\ang{90},\ 0.40)$ の行は $(0,\ -1,\ 0)$、$(1,\ 0,\ 0.40)$、$(0,\ 0,\ 1)$。$A(\ang{-90},\ 0.30)$ の右の列は $(0,\ -0.30,\ 1)$。積の右の列は $(0.30,\ 0.40,\ 1)$ で、手は $(0.30,\ 0.40)$\,m}
\end{jibun}

\begin{phreidai}[関節を回したときの手の速さ]\label{ph:力-7:ex8}
\bkreidairef{5}の腕（$\theta_1=\ang{30}$、$\theta_2=\ang{45}$）で、肩を\qty{0.50}{rad/s}、ひじを\qty{1.0}{rad/s}（どちらも左回り）で回しています。この瞬間の手の速度と速さを求めましょう。
\phbamen
2つの関節を同時に回すと、手はどの向きにどれだけの速さで動くかを調べます。
\phmushi
腕のたわみを無視します。角速度は、この瞬間の値です。
\phzu
$\dot\theta_1=0.50$、$\dot\theta_2=1.0$（rad/s）、$\theta_{12}=\ang{75}$。
\phriyu
関節の角速度から手の速度を出すので、ヤコビ行列の式（⑦）を使います。
\phtateru
\[J=\begin{pmatrix}-0.30\sin\ang{30}-0.25\sin\ang{75}&-0.25\sin\ang{75}\\0.30\cos\ang{30}+0.25\cos\ang{75}&0.25\cos\ang{75}\end{pmatrix}\]
計算すると\[J=\begin{pmatrix}-0.391&-0.241\\0.325&0.065\end{pmatrix}\]
\phtoku
\[\dot x=-0.391\times0.50-0.241\times1.0\approx-0.437\]
\[\dot y=0.325\times0.50+0.065\times1.0\approx0.227\]
速さ $\sqrt{0.437^2+0.227^2}\approx\qty{0.49}{m/s}$。
\phtashikame
単位：m×rad/s＝m/s。\\
1列めの確かめ：$J$ の1列めは、手の位置 $(0.325,\ 0.391)$ を\ang{90}回した $(-0.391,\ 0.325)$ になっています。\\
極端な場合：ひじだけを止める（$\dot\theta_2=0$）と、手は肩を中心とする円の接線の向きに $0.509\times0.50\approx\qty{0.25}{m/s}$ で動きます。
\kotae 速度 約 $(-0.44,\ 0.23)$\,m/s、速さ 約\qty{0.49}{m/s}
\end{phreidai}
\begin{jibun}
$L_1=\qty{0.40}{m}$、$L_2=\qty{0.30}{m}$、$\theta_1=\ang{20}$、$\theta_2=\ang{50}$ の腕で、肩を\qty{0.40}{rad/s}、ひじを\qty{0.50}{rad/s}で回すときの手の速度と速さを求めましょう。
\bkjans{$J$ の1行めは $(-0.419,\ -0.282)$、2行めは $(0.478,\ 0.103)$。速度 約 $(-0.31,\ 0.24)$\,m/s、速さ 約\qty{0.39}{m/s}}
\end{jibun}

\begin{phreidai}[エンジンのピストンの速さ]\label{ph:力-7:ex9}
クランクの半径\qty{0.040}{m}（行程\qty{0.080}{m}）、連接棒の長さ\qty{0.14}{m}のエンジンが、毎分3000回転しています。（1）ピストンの平均の速さ、（2）クランクが\ang{90}のときのピストンの位置（クランクの軸からの距離）を求めましょう。
\phbamen
クランクが一定の速さで回り、ピストンが行ったり来たりします。
\phmushi
部品の変形と、すきまを無視します。回る速さは一定とします。
\phzu
⑥の図。$r=0.040$、$l=0.14$（m）、毎秒 $n$ 回転。
\phriyu
1回転で行程を2回動くので、平均の速さは $4rn$（⑥）。位置は⑥の $x=r\cos\theta+\sqrt{l^2-r^2\sin^2\theta}$。
\phtateru
\[n=\frac{3000}{60}=50\ \text{回転/s},\qquad \text{平均の速さ}=4\times0.040\times50\]
\[x=0.040\cos\ang{90}+\sqrt{0.14^2-0.040^2}\]
\phtoku
平均の速さ $=\qty{8.0}{m/s}$。
\[x=0+\sqrt{0.0196-0.0016}=\sqrt{0.0180}\approx\qty{0.134}{m}\]
\phtashikame
単位：m×1/s＝m/s。\\
大きさ：自動車のエンジンのピストンの平均の速さは、ふつう秒速数m〜十数mです。\\
位置：行程のまん中は $\dfrac{0.18+0.10}{2}=0.14$ で、\ang{90}のピストンはそれより\qty{6}{mm}ほどクランク側にいます（⑥のグラフ）。
\kotae （1）\qty{8.0}{m/s}（2）約\qty{0.134}{m}
\end{phreidai}
\begin{jibun}
クランクの半径\qty{0.035}{m}、連接棒の長さ\qty{0.12}{m}のエンジンが、毎分6000回転しています。ピストンの平均の速さと、クランクが\ang{90}のときのピストンの位置を求めましょう。
\bkjans{平均 $4\times0.035\times100=\qty{14}{m/s}$。位置 $\sqrt{0.12^2-0.035^2}\approx\qty{0.115}{m}$}
\end{jibun}

\begin{phreidai}[4節リンクのてこの角を表計算で]\label{ph:力-7:ex10}
⑥の4節リンク（クランク $a=\qty{2}{cm}$、つなぎ $b=\qty{5}{cm}$、てこ $c=\qty{4}{cm}$、固定 $d=\qty{6}{cm}$）で、クランクの角 $\varphi=\ang{90}$ のときの、てこの角 $\psi$ を求めましょう。また、表計算で $\varphi$ を\ang{30}ずつ変えて、$\psi$ の変わり方を調べましょう。
\phbamen
クランク $\mathrm{A_0A}$ を回すと、てこ $\mathrm{B_0B}$ が行ったり来たりします。てこの角を、クランクの角で表したい。
\phmushi
棒の太さと、ピンのすきまを無視します。
\phzu
$\mathrm{A_0}=(0,\ 0)$、$\mathrm{B_0}=(6,\ 0)$、$\mathrm{A}=(2\cos\varphi,\ 2\sin\varphi)$。$e=\mathrm{AB_0}$、$\mathrm{B_0}$ からAを見た向きの角 $\delta$、三角形 $\mathrm{AB_0B}$ の $\mathrm{B_0}$ の角 $\gamma$。てこが上にある形で、$\psi=\delta-\gamma$。
\phriyu
三角形 $\mathrm{AB_0B}$ の3辺 $e$、$b$、$c$ がわかれば、余弦定理で $\gamma$ が出ます（逆運動学と同じ考え方）。
\phtateru
\[e=\sqrt{(2\cos\varphi-6)^2+(2\sin\varphi)^2},\qquad \cos\gamma=\frac{c^2+e^2-b^2}{2ce}\]
\phtoku
$\varphi=\ang{90}$ なら $\mathrm{A}=(0,\ 2)$、$e=\sqrt{36+4}\approx6.32$。$\mathrm{B_0}$ からAへの向きは、左へ6、上へ2なので、$\delta=\ang{180}-\tan^{-1}\dfrac{2}{6}\approx\ang{161.6}$。
\[\cos\gamma=\frac{16+40-25}{2\times4\times6.32}\approx0.613,\qquad \gamma\approx\ang{52.2}\]
\[\psi=\ang{161.6}-\ang{52.2}=\ang{109.4}\]
表計算では、A列に $\varphi$ を入れ、Aの点の座標、$e$、$\delta$（\texttt{ATAN2} で求め、負になったら\ang{360}を足す）、$\gamma$（\texttt{ACOS}）、$\psi=\delta-\gamma$ を順に計算させます。$\psi$ が\ang{360}以上なら\ang{360}を引きます。$\varphi$ を\ang{30}ずつ変えると、次のようになります。
\begin{bkcenter}
\begin{bktable}{colspec={Q[c,m,wd=11mm] Q[c,m,wd=11mm] Q[c,m,wd=11mm] Q[c,m,wd=11mm] Q[c,m,wd=11mm] Q[c,m,wd=11mm] Q[c,m,wd=11mm]}, rowsep=1pt, cells={font=\footnotesize}}
$\varphi$ & 0 & 30 & 60 & 90 & 120 & 150 \\
$\psi$ & 102.6 & 93.7 & 97.6 & 109.3 & 124.3 & 138.8 \\
$\varphi$ & 180 & 210 & 240 & 270 & 300 & 330 \\
$\psi$ & 149.2 & 153.5 & 152.1 & 146.2 & 135.8 & 120.1 \\
\end{bktable}
\end{bkcenter}
\phtashikame
角の単位は°。表の\ang{90}の値109.3は、手計算（109.4。途中で丸めたため）と合います。\\
$\psi$ は約\ang{94}〜約\ang{154}の間で行ったり来たりし、クランクが1回転すると1往復します。
\kotae $\psi\approx\ang{109}$。てこは約\ang{94}〜約\ang{154}（約\ang{60}のはば）を往復する
\end{phreidai}
\begin{jibun}
同じ4節リンクで、$\varphi=\ang{180}$ のときの $e$、$\gamma$、$\psi$ を求めましょう。
\bkjans{$\mathrm{A}=(-2,\ 0)$、$e=8$、$\delta=\ang{180}$。$\cos\gamma=\dfrac{16+64-25}{64}\approx0.859$、$\gamma\approx\ang{30.8}$、$\psi\approx\ang{149.2}$}
\end{jibun}

\begin{phreidai}[歯車列の速さとトルク]\label{ph:力-7:ex11}
モータ（毎分3000回転、トルク\qty{0.50}{N.m}）の軸に20歯の歯車をつけ、60歯の歯車とかみ合わせます。60歯の歯車と同じ軸に15歯の歯車をつけ、それを45歯の歯車とかみ合わせます。（1）60歯の軸、（2）45歯の軸の、回る速さとトルクを求めましょう。
\phbamen
2組の歯車で、モータの速さを落とし、トルクを大きくする減速機です。
\phmushi
歯車の摩擦などの損失を無視します（効率1）。軸と歯車の慣性モーメントも無視します（一定の速さで回っているとき）。
\phzu
1組め $z_1=20$、$z_2=60$。2組め $z_3=15$、$z_4=45$。
\phriyu
回る速さの比は歯の数の逆の比、トルクの比は歯の数の比です（⑨）。2組は、かけ算でつなぎます。
\phtateru
\[n_2=3000\times\frac{20}{60},\qquad N_2=0.50\times\frac{60}{20}\]
\[n_4=n_2\times\frac{15}{45},\qquad N_4=N_2\times\frac{45}{15}\]
\phtoku
（1）毎分1000回転、\qty{1.5}{N.m}。（2）毎分約333回転、\qty{4.5}{N.m}。
\phtashikame
動力（$P=N\omega$）：モータと45歯の軸で、
\[0.50\times\frac{3000\times2\pi}{60}\approx\qty{157}{W},\qquad 4.5\times\frac{333.3\times2\pi}{60}\approx\qty{157}{W}\]
同じです。\\
回る向き：2回かみ合うので、45歯の軸はモータと同じ向きです。\\
効率：1組の効率を0.95とすると、2組で $0.95^2\approx0.90$、トルクは約\qty{4.1}{N.m}になります。
\kotae （1）毎分1000回転・\qty{1.5}{N.m}（2）毎分約330回転・\qty{4.5}{N.m}
\end{phreidai}
\begin{jibun}
毎分1800回転・\qty{0.20}{N.m}のモータに12歯の歯車をつけ、48歯とかみ合わせ、その軸の15歯を45歯とかみ合わせました。損失がないとして、最後の軸の回る速さとトルクを求めましょう。
\bkjans{減速比 $4\times3=12$。毎分150回転、\qty{2.4}{N.m}}
\end{jibun}

\begin{phreidai}[自転車のギヤとスピード]\label{ph:力-7:ex12}
前の歯車（44歯）と後ろの歯車（16歯）をチェーンでつないだ自転車で、ペダルを毎分60回転でこぎます。後ろの車輪の半径を\qty{0.33}{m}として、自転車の速さを求めましょう。
\phbamen
ペダルの回転が、チェーンで後ろの歯車に伝わり、後ろの車輪が道を転がります。
\phmushi
チェーンのゆるみ、タイヤのすべりを無視します。後ろの歯車と後ろの車輪は同じ軸で、同じ速さで回ります。
\phzu
ペダル（前の歯車）の角速度 $\omega_1$、後ろの車輪の角速度 $\omega_2$、半径 $r=\qty{0.33}{m}$。
\phriyu
チェーンがすべらないので、ベルトと同じく $\dfrac{\omega_2}{\omega_1}=\dfrac{z_1}{z_2}$（⑨）。車輪は転がるので $v=r\omega_2$（①）。
\phtateru
\[\omega_1=\frac{60\times2\pi}{60}=2\pi\ \text{rad/s},\qquad \omega_2=2\pi\times\frac{44}{16},\qquad v=0.33\,\omega_2\]
\phtoku
\[\omega_2=2\pi\times2.75\approx\qty{17.3}{rad/s},\qquad v=0.33\times17.3\approx\qty{5.7}{m/s}\]
\phtashikame
\qty{5.7}{m/s}は約\qty{21}{km/h}で、ふつうに走る自転車の速さです。\\
2倍にすると：後ろを22歯にすると、$\dfrac{44}{22}=2$ で約\qty{4.1}{m/s}。歯の多い後ろの歯車は、坂道で使う「軽い」ギヤです（速さが出ないかわりに、こぐ力が小さくてすむ）。
\kotae 約\qty{5.7}{m/s}（約\qty{21}{km/h}）
\end{phreidai}
\begin{jibun}
前48歯・後ろ12歯で、ペダルを毎分80回転でこぎます。後ろの車輪の半径\qty{0.33}{m}として、自転車の速さを求めましょう。
\bkjans{$\omega_2=\dfrac{80\times2\pi}{60}\times\dfrac{48}{12}\approx\qty{33.5}{rad/s}$、$v=0.33\times33.5\approx\qty{11}{m/s}$（約\qty{40}{km/h}）}
\end{jibun}

\begin{phreidai}[回転台の上の物がすべり出す回転数]\label{ph:力-7:ex13}
回転台の中心から\qty{0.20}{m}の所に、小さな物をのせました。物と台の間の静止摩擦係数を0.40とします。台の回転を少しずつ速くすると、毎分何回転で物はすべり出しますか。
\phbamen
台がゆっくり加速するので、物はほぼ等速円運動をしています。速くなると、摩擦が足りなくなり、すべり出します。
\phmushi
物の大きさと、台の加速の影響を無視します。空気の抵抗も無視します。
\phzu
\noindent\begin{minipage}[c]{0.36\linewidth}\centering
\begin{tikzpicture}[x=1mm, y=1mm]
  \filldraw[fill=bkTheme!8, draw=bkTheme!70, line width=0.8pt] (0,0) circle (16);
  \fill[bkInk] (0,0) circle (0.8);
  \draw[bkInk, line width=0.7pt, ->, >={Stealth[length=3pt]}] (60:18.5) arc (60:100:18.5);
  \node[font=\sffamily\footnotesize, anchor=south] at (80:19) {$\omega$};
  \filldraw[fill=bkOrange!60, draw=bkInk, line width=0.5pt] (8.5,-1.5) rectangle (11.5,1.5);
  \draw[bk arrow, bkBlue, line width=1.1pt] (8.5,0) -- (2,0);
  \draw[bk arrow, bkRed, line width=1.1pt] (11.5,0) -- (19,0);
  \node[font=\sffamily\footnotesize, text=bkBlue!80!black, anchor=north] at (4.5,-0.8) {$f$};
  \node[font=\sffamily\footnotesize, text=bkRed!80!black, anchor=north] at (16,-0.8) {$mr\omega^2$};
\end{tikzpicture}
\end{minipage}\hfill
\begin{minipage}[c]{0.6\linewidth}
上から見た図。台といっしょに回る座標で考えます。物にかかる水平な力は、静止摩擦力 $f$（中心向き）と遠心力 $mr\omega^2$（外向き）。上下には、重力 $mg$ と垂直抗力 $N=mg$ がつり合っています。
\end{minipage}
\phriyu
台の上から見ると、物は止まっているので、遠心力と静止摩擦力がつり合います（⑧）。摩擦力の最大は $\mu N$（力-2）。
\phtateru
すべり出す直前：
\[mr\omega^2=\mu mg\qquad\omega=\sqrt{\frac{\mu g}{r}}\]
\phtoku
\[\omega=\sqrt{\frac{0.40\times9.8}{0.20}}=\sqrt{19.6}\approx\qty{4.43}{rad/s}\]
\[\text{毎分}\ \frac{4.43\times60}{2\pi}\approx42\ \text{回転}\]
\phtashikame
地面から見る別の方法：物は等速円運動をしていて、中心向きの力（向心力）$mr\omega^2$ を、静止摩擦力が出しています（力-2）。式は同じになります。\\
大きさ：レコードプレーヤの毎分33回転や45回転くらいです。\\
極端な場合：$m$ が式から消えるので、重い物も軽い物も同じ回転数ですべり出します。中心に近い（$r$ が小さい）ほど、すべりにくくなります。
\kotae 毎分約42回転（約\qty{4.4}{rad/s}）
\end{phreidai}
\begin{jibun}
中心から\qty{0.10}{m}の所にのせた物（静止摩擦係数0.50）が、すべり出すのは毎分何回転のときですか。
\bkjans{$\omega=\sqrt{\dfrac{0.50\times9.8}{0.10}}=\qty{7.0}{rad/s}$、毎分約67回転}
\end{jibun}

\begin{phreidai}[カーブでトラックが横転する速さ]\label{ph:力-7:ex14}
重心の高さ\qty{1.5}{m}、左右の車輪の間の半分が\qty{1.0}{m}のトラックが、半径\qty{50}{m}の平らなカーブを曲がります。（1）外側の車輪のまわりに倒れ始める速さを求めましょう。（2）タイヤと道の静止摩擦係数を0.7とすると、すべり出すのとどちらが先ですか。
\phbamen
カーブを曲がるトラックの中の人には、外へ引っぱられる力（遠心力）を感じます。速すぎると、倒れるか、すべります。
\phmushi
サスペンション（ばね）でかたむくこと、荷のずれ、道のかたむき（カント）を無視します。トラックは剛体で、重心は左右のまん中とします。
\phzu
\noindent\begin{minipage}[c]{0.42\linewidth}\centering
\begin{tikzpicture}[x=1mm, y=1mm]
  \draw[bkInk, line width=0.9pt] (-18,0) -- (22,0);
  \filldraw[fill=bkBlue!15, draw=bkInk, line width=0.8pt] (-12,3) rectangle (12,26);
  \fill[bkInk] (-9,0) ++(0,2.5) circle (2.5) (9,0) ++(0,2.5) circle (2.5);
  \fill[bkInk] (0,16) circle (1);
  \node[font=\sffamily\footnotesize, anchor=south east] at (-0.6,16.4) {G};
  \draw[bk arrow, bkRed, line width=1.2pt] (0,16) -- (0,6);
  \node[font=\sffamily\footnotesize, text=bkRed!80!black, anchor=east] at (-0.6,8) {$mg$};
  \draw[bk arrow, bkOrange!90!black, line width=1.2pt] (0,16) -- (13,16);
  \node[font=\sffamily\footnotesize, text=bkOrange!80!black, anchor=west] at (13.6,16) {$\dfrac{mv^2}{R}$};
  \fill[bkGreen!60!black] (9,0) circle (1);
  \node[font=\sffamily\footnotesize, anchor=north] at (9,-0.8) {外側の車輪};
  \draw[bkGray, line width=0.5pt, {Stealth[length=3pt]}-{Stealth[length=3pt]}] (-16,0) -- (-16,16);
  \node[font=\sffamily\footnotesize, anchor=east] at (-16.5,8) {$h$};
  \draw[bkGray, line width=0.5pt, {Stealth[length=3pt]}-{Stealth[length=3pt]}] (0,-6) -- (9,-6);
  \node[font=\sffamily\footnotesize, anchor=north] at (4.5,-6.4) {$b$};
\end{tikzpicture}
\end{minipage}\hfill
\begin{minipage}[c]{0.56\linewidth}
後ろから見た図（カーブの中心は左）。トラックといっしょに動く座標で考え、重心Gに外向きの遠心力 $\dfrac{mv^2}{R}$ をかけます。$h=\qty{1.5}{m}$、$b=\qty{1.0}{m}$、$R=\qty{50}{m}$。
\end{minipage}
\phriyu
倒れ始めるのは、外側の車輪のまわりで、遠心力が倒す向きに回すモーメントが、重力がもどす向きに回すモーメントをこえるときです（力-5 のつり合い）。内側の車輪が浮く直前は、内側の車輪の力が0です。
\phtateru
（1）外側の車輪のまわりのモーメントのつり合い：
\[\frac{mv^2}{R}\times h=mg\times b\qquad v=\sqrt{\frac{gbR}{h}}\]
（2）すべり出す速さ：$\dfrac{mv^2}{R}=\mu mg$ より $v=\sqrt{\mu gR}$。
\phtoku
（1）$v=\sqrt{\dfrac{9.8\times1.0\times50}{1.5}}=\sqrt{326.7}\approx\qty{18.1}{m/s}$（約\qty{65}{km/h}）\\
（2）$v=\sqrt{0.7\times9.8\times50}=\sqrt{343}\approx\qty{18.5}{m/s}$（約\qty{67}{km/h}）。倒れるほうが先です。
\phtashikame
単位：$\sqrt{\text{m/s}^2\times\text{m}\times\text{m}\div\text{m}}=\text{m/s}$。\\
くらべ方：$\dfrac{b}{h}<\mu$ なら倒れるほうが先。ここでは $\dfrac{1.0}{1.5}=0.67<0.7$。\\
実際の車はサスペンションでかたむき、重心が外へずれるので、この計算より低い速さで倒れることがあります。
\kotae （1）約\qty{18}{m/s}（約\qty{65}{km/h}）（2）倒れるほうが先
\end{phreidai}
\begin{jibun}
重心の高さ\qty{0.55}{m}、左右の車輪の間の半分が\qty{0.75}{m}の乗用車が、半径\qty{50}{m}の平らなカーブを曲がります。倒れ始める速さと、すべり出す速さ（$\mu=0.7$）を求め、どちらが先か答えましょう。
\bkjans{倒れる $\sqrt{\dfrac{9.8\times0.75\times50}{0.55}}\approx\qty{26}{m/s}$、すべる $\sqrt{0.7\times9.8\times50}\approx\qty{18.5}{m/s}$。すべるほうが先}
\end{jibun}

\begin{phreidai}[地球のコリオリの力：上空の風]\label{ph:力-7:ex15}
北緯\ang{35}の上空を、ジェット気流が東へ\qty{100}{m/s}でふいています。地球の自転によるコリオリの力が、空気に与える水平な加速度の大きさを求めましょう。向きはどちらですか。
\phbamen
地球といっしょに回る私たちから見た、空気の動きです。
\phmushi
地球の形（完全な球でないこと）、空気の上下の動きを無視します。地球の自転の角速度は一定とします。
\phzu
地球の角速度 $\omega=\qty{7.29e-5}{rad/s}$、緯度 $\phi=\ang{35}$、風の速さ $v=\qty{100}{m/s}$。
\phriyu
緯度 $\phi$ の地面は、真上の軸のまわりに $\omega\sin\phi$ で回っているので、水平なコリオリの力は $2m\omega v\sin\phi$（⑧）。加速度は、これを $m$ でわったものです。
\phtateru
\[a=2\omega v\sin\phi=2\times7.29\times10^{-5}\times100\times\sin\ang{35}\]
\phtoku
\[a\approx2\times7.29\times10^{-5}\times100\times0.574\approx\qty{0.0084}{m/s^2}\]
向きは、北半球なので、進む向き（東）の右、つまり南向きです。
\phtashikame
単位：rad/s×m/s＝m/s$^2$。\\
大きさ：$g$ の約\num{0.00085}倍。とても小さい力ですが、大きな空気のかたまりが何時間も動くと、向きを大きく曲げます。\\
極端な場合：赤道（$\phi=0$）では水平なコリオリの力は0です。
\kotae 約\qty{0.0084}{m/s^2}、南向き（進む向きの右）
\end{phreidai}
\begin{jibun}
北緯\ang{45}で、\qty{50}{m/s}で北へ進む空気にはたらく、水平なコリオリの力による加速度の大きさと向きを求めましょう。
\bkjans{$2\times7.29\times10^{-5}\times50\times0.707$ で約\qty{0.0052}{m/s^2}。東向き（進む向きの右）}
\end{jibun}

\begin{phreidai}[ジャイロ：つるした車輪の歳差運動]\label{ph:力-7:ex16}
質量\qty{2.0}{kg}、軸のまわりの慣性モーメント\qty{0.10}{kg.m^2}の自転車の車輪を\qty{30}{rad/s}で回し、軸の片方のはしを、車輪の中心から\qty{0.20}{m}の所でひもでつるしました。軸が水平にひものまわりを回る角速度と、1周する時間を求めましょう。
\phbamen
車輪は落ちずに、軸が水平のまま、ゆっくりひものまわりを回ります（⑩）。
\phmushi
軸の質量、空気の抵抗、軸受けの摩擦を無視します。車輪は十分速く回っているとします。
\phzu
$m=\qty{2.0}{kg}$、$d=\qty{0.20}{m}$、$I=\qty{0.10}{kg.m^2}$、$\omega=\qty{30}{rad/s}$。
\phriyu
重力のトルクが角運動量の向きを変える歳差運動の式 $\Omega=\dfrac{mgd}{I\omega}$（⑩）を使います。
\phtateru
\[\Omega=\frac{2.0\times9.8\times0.20}{0.10\times30},\qquad T=\frac{2\pi}{\Omega}\]
\phtoku
\[\Omega=\frac{3.92}{3.0}\approx\qty{1.3}{rad/s},\qquad T=\frac{2\pi}{1.31}\approx\qty{4.8}{s}\]
\phtashikame
単位：N·m÷（kg·m$^2$×rad/s）＝1/s。\\
近似の確かめ：$\Omega=1.3$ は $\omega=30$ よりずっと小さいので、⑩の式が使えます。\\
2倍にすると：$\omega$ を\qty{60}{rad/s}にすると、$\Omega$ は半分の約\qty{0.65}{rad/s}。
\kotae 約\qty{1.3}{rad/s}、約\qty{4.8}{s}で1周
\end{phreidai}
\begin{jibun}
質量\qty{1.5}{kg}、慣性モーメント\qty{0.050}{kg.m^2}の車輪を\qty{40}{rad/s}で回し、中心から\qty{0.15}{m}の所でつるしました。歳差運動の角速度を求めましょう。
\bkjans{$\Omega=\dfrac{1.5\times9.8\times0.15}{0.050\times40}\approx\qty{1.1}{rad/s}$}
\end{jibun}

%% ================================================================
\bksection{tsumazuki}{まちがえやすいところをチェック}
\begin{tsumazuki}{転がる車輪の上の点の速さ}
\bkng{車輪のどの点も、自転車と同じ速さ $v$ で進んでいる。}
\bkok{いちばん上は $2v$、中心は $v$、地面にふれた点は0。}
進む速さに、回る速さ $r\omega$ がたされたり引かれたりするからです（①）。
\end{tsumazuki}
\begin{tsumazuki}{転がる物の摩擦}
\bkng{斜面を転がり下りる球の摩擦力は、エネルギーを熱にしている。}
\bkok{すべらずに転がるとき、静止摩擦力は仕事をしない。}
力がかかる点（ふれた点）が、その瞬間に止まっているからです。摩擦力は、エネルギーを回転のエネルギーに分けているだけです（⑤）。
\end{tsumazuki}
\begin{tsumazuki}{逆運動学の答えの数}
\bkng{手の位置を決めれば、関節の角は1つに決まる。}
\bkok{ふつう2通り（ひじが上と下）。届かない所では0通り。}
余弦定理で求めた $\cos\theta_2$ に、＋と－の2つの角があるからです（②・⑦）。
\end{tsumazuki}
\begin{tsumazuki}{歯車で強くなる？}
\bkng{歯車を使うと、速さもトルクも大きくできる。}
\bkok{速さを落とせばトルクが大きくなり、速さを上げればトルクは小さくなる。}
損失がなければ、動力（速さ×トルク）は変わらないからです（⑨）。
\end{tsumazuki}
\begin{tsumazuki}{遠心力の正体}
\bkng{遠心力は、回る物を外へ引っぱる本当の力だ。}
\bkok{遠心力は、回る座標の中から見るときだけに加える、見かけの力（慣性力）。}
地面から見ると、外へ引く力はありません。中心向きの力が足りないと、物はまっすぐ進もうとして、外へ出ていくのです（⑧）。
\end{tsumazuki}
\begin{tsumazuki}{ふろの栓のうず}
\bkng{北半球では、ふろの栓をぬくと、コリオリの力でかならず左回りのうずになる。}
\bkok{ふつうのふろでは、うずの向きは、水に残っていた動きや形で決まる。}
ふろの大きさでは、コリオリの力はとても小さいからです。1962年にアメリカで、直径約\qty{1.8}{m}の丸い水そうの水を24時間静かにしてから栓をぬく実験をすると、やっと左回りのうずが確かめられました。
\end{tsumazuki}

%% ================================================================
\bksection{phgear}{工学ではこう使う}
\begin{phkougaku}[水平多関節ロボット（スカラロボット）]
電子部品を基板にのせる工場などでは、\textbf{水平多関節ロボット}（スカラロボット）が働いています。この章の2リンクの腕を水平に置き、手の先に上下に動く軸と回る軸をつけたものです。肩とひじの角を逆運動学で計算し、部品を決まった位置へ運びます。
\begin{bkcenter}
\begin{tikzpicture}[x=1mm, y=1mm]
  \fill[bkGray!30] (-8,0) rectangle (8,4);
  \fill[bkGray!55] (-5,4) rectangle (5,30);
  \fill[bkBlue!60] (-5,26) rectangle (38,34);
  \fill[bkBlue!40] (32,30) rectangle (70,37);
  \fill[bkGray!70] (64,10) rectangle (67,40);
  \filldraw[fill=white, draw=bkInk, line width=0.6pt] (0,30) circle (2.4) (35,33.5) circle (2.2);
  \fill[bkOrange!70] (60,4) rectangle (71,8);
  \draw[bk arrow, bkRed, line width=1pt] (74,30) -- (74,16);
  \node[font=\sffamily\footnotesize, anchor=west] at (75,23) {上下};
  \node[font=\sffamily\footnotesize, anchor=south] at (0,38) {肩};
  \node[font=\sffamily\footnotesize, anchor=south] at (35,40) {ひじ};
  \node[font=\sffamily\footnotesize, anchor=north] at (65.5,3.4) {部品};
\end{tikzpicture}
\bkfigcaption{スカラロボットのしくみ（横から見た図）。肩とひじは、上から見ると、この章の2リンクの腕と同じに動く。}
\end{bkcenter}
\begin{phketa}[ロボットと回る物の数]
\phketaline{スカラロボットの位置くり返し精度（同じ場所にもどる正確さ）}{約±\qty{0.02}{mm}}
\phketaline{垂直多関節ロボット（6つの関節）の位置くり返し精度}{約±\num{0.02}〜\qty{0.03}{mm}}
\phketaline{地球の自転の角速度}{約\qty{7.29e-5}{rad/s}}
\phketaline{洗濯機の脱水（毎分800回転、半径\qty{0.25}{m}）の $r\omega^2$}{約\qty{1800}{m/s^2}}
\phboxhead{答えが桁で合っているか}
腕の長さが\qty{0.5}{m}のとき、\qty{0.02}{mm}は腕の長さの約 $\dfrac{1}{\num{25000}}$ です。角にすると約 $4\times10^{-5}$\,rad（$s=r\theta$）で、関節のモータの角は、とても細かく決めなければなりません。答えの角を\ang{0.1}の位で丸めると、手の位置は約 $0.5\times0.0017\approx\qty{0.9}{mm}$ ずれます。
\end{phketa}
\phboxhead{減速機：モータとタイヤをつなぐ}
電気自動車や電車のモータは、速く回るほうが小さく軽く作れます。そこで、モータとタイヤの間に歯車の減速機を入れ、回る速さを落として、タイヤを回すトルクを大きくしています（⑨）。ロボットの関節にも、小さくて減速比の大きい減速機が入っています。
\phboxhead{スマホのジャイロセンサ}
スマホの向きを知るセンサの1つ（ジャイロセンサ）は、小さなおもりを細かくふるわせておき、スマホが回るとおもりにはたらくコリオリの力をはかって、回る速さを求めています（⑧）。手のひらにのる電子部品の中で、この章の「見かけの力」が働いているのです。
\end{phkougaku}
\phstation{機構学,ロボット工学,機械力学,車両の運動}{平面運動、瞬間中心、順運動学、逆運動学、同次変換、ヤコビ行列、特異姿勢、減速比、コリオリの力、ジャイロ効果}

\begin{column}{地球が回っていることを、振り子で見せた人}
1851年、フランスのフーコーは、パリで、長いひもでつるした重い振り子をゆらす実験を人々に見せました。振り子がゆれる面は、何時間もかけて、少しずつ右回りに回っていきました。振り子のゆれる面が回ったのではなく、振り子の下の地面（地球）が回っていたのです。回る部屋から見ると、まっすぐゆれるはずのおもりに、コリオリの力がはたらいて、ゆれる面が回って見えます。ゆれる面が1周する時間は、1恒星日÷$\sin$（緯度）で、パリでは約32時間、東京では約41時間です。日本の科学館にも、フーコーの振り子が置かれている所があります。
\end{column}

%% ================================================================
\bkendproblems
\bkproblemgroup{基本}{4}
\bkproblem 半径\qty{0.35}{m}の車輪が、すべらずに\qty{14}{m/s}で転がっています。車輪の角速度と、いちばん上の点の速さを求めましょう。
\bkpans{$\omega=\dfrac{14}{0.35}=\qty{40}{rad/s}$、上の点\qty{28}{m/s}}
\bkproblem 角\ang{15}、長さ\qty{4.0}{m}の斜面の上から、一様な球をそっと放しました。下に着くまでの時間と、下での速さを求めましょう（$\sin\ang{15}=0.259$）。
\bkpans{$a=\dfrac{9.8\times0.259}{1.4}\approx\qty{1.81}{m/s^2}$、$t=\sqrt{\dfrac{8.0}{1.81}}\approx\qty{2.1}{s}$、$v=at\approx\qty{3.8}{m/s}$}
\bkproblem $L_1=\qty{0.20}{m}$、$L_2=\qty{0.15}{m}$ の腕で、$\theta_1=\ang{60}$、$\theta_2=\ang{30}$ のときの手の位置を求めましょう。
\bkpans{$x=0.20\cos\ang{60}+0.15\cos\ang{90}=0.10$、$y=0.20\sin\ang{60}+0.15\sin\ang{90}\approx0.32$。$(0.10,\ 0.32)$\,m}
\bkproblem 毎分1200回転・\qty{0.80}{N.m}のモータに15歯の歯車をつけ、45歯の歯車とかみ合わせました。損失がないとして、45歯の軸の回る速さとトルクを求めましょう。
\bkpans{毎分400回転、\qty{2.4}{N.m}}
\bkproblem 前40歯・後ろ20歯の自転車で、ペダルを毎分60回転でこぎます。後ろの車輪の半径を\qty{0.33}{m}として、自転車の速さを求めましょう。
\bkpans{$\omega_2=2\pi\times2\approx\qty{12.6}{rad/s}$、$v\approx\qty{4.1}{m/s}$（約\qty{15}{km/h}）}
\bkproblem 回転台の中心から\qty{0.15}{m}の所にのせた物（静止摩擦係数0.30）は、毎分何回転ですべり出しますか。
\bkpans{$\omega=\sqrt{\dfrac{0.30\times9.8}{0.15}}\approx\qty{4.4}{rad/s}$、毎分約42回転}
\bkproblem クランクの半径\qty{0.045}{m}のエンジンが毎分2000回転しています。ピストンの平均の速さを求めましょう。
\bkpans{$4\times0.045\times\dfrac{2000}{60}=\qty{6.0}{m/s}$}
\bkproblemgroup{チャレンジ}{5}
\bkproblem $L_1=\qty{0.30}{m}$、$L_2=\qty{0.25}{m}$ の腕で、手を点 $(0.30,\ 0.30)$（m）に置く角を2通り求めましょう。
\bkpans{$d\approx0.424$、$\cos\angle\mathrm{OEP}\approx-0.183$、$\theta_2=\pm\ang{79.4}$。$\alpha=\ang{45}$、$\beta\approx\ang{35.4}$。$(\ang{9.6},\ \ang{79.4})$ と $(\ang{80.4},\ \ang{-79.4})$}
\bkproblem 北緯\ang{35}で、\qty{20}{m/s}で走る車にはたらく、水平なコリオリの力による加速度の大きさを求め、$g$ の何倍か答えましょう。
\bkpans{$2\times7.29\times10^{-5}\times20\times0.574$ で約\qty{0.0017}{m/s^2}。$g$ の約\num{0.00017}倍}
\bkproblem 質量\qty{1.0}{kg}、慣性モーメント\qty{0.020}{kg.m^2}の車輪を\qty{50}{rad/s}で回し、中心から\qty{0.10}{m}の所でつるしました。歳差運動の角速度を求めましょう。
\bkpans{$\Omega=\dfrac{1.0\times9.8\times0.10}{0.020\times50}=\qty{0.98}{rad/s}$}
\bkproblem 重心の高さ\qty{1.2}{m}、左右の車輪の間の半分が\qty{1.0}{m}のバスが、半径\qty{40}{m}の平らなカーブを曲がります。静止摩擦係数を0.6として、倒れ始める速さとすべり出す速さを求め、どちらが先か答えましょう。
\bkpans{倒れる $\sqrt{\dfrac{9.8\times1.0\times40}{1.2}}\approx\qty{18}{m/s}$、すべる $\sqrt{0.6\times9.8\times40}\approx\qty{15}{m/s}$。すべるほうが先}
\bkproblem $L_1=\qty{0.30}{m}$、$L_2=\qty{0.25}{m}$ の腕が、$\theta_1=\ang{0}$、$\theta_2=\ang{90}$ の姿勢で、肩だけを\qty{1.0}{rad/s}で回しています。手の速度と速さを求めましょう。
\bkpans{手は $(0.30,\ 0.25)$。$J$ の1列めは $(-0.25,\ 0.30)$ なので、速度 $(-0.25,\ 0.30)$\,m/s、速さ約\qty{0.39}{m/s}}
\bkproblem ⑥の4節リンクで、$\varphi=\ang{0}$ のときのてこの角 $\psi$ を求めましょう。
\bkpans{$\mathrm{A}=(2,\ 0)$、$e=4$、$\delta=\ang{180}$、$\cos\gamma=\dfrac{16+16-25}{32}\approx0.219$、$\gamma\approx\ang{77.4}$、$\psi\approx\ang{102.6}$}
\phnyushi
\bkproblem 質量 $m$、半径 $r$ の一様な円柱（$I=\dfrac{1}{2}mr^2$）を、回さずに速さ $v_0$ で水平な床にすべらせました。動摩擦係数を $\mu'$ とします。（1）すべらずに転がり始めるときの速さ、（2）それまでの時間、（3）失われた運動エネルギーは、はじめの運動エネルギーの何倍ですか。
\bkpans{（1）$v=\dfrac{2}{3}v_0$（2）$t=\dfrac{v_0}{3\mu'g}$（3）$\dfrac{1}{3}$ 倍}
\bkproblem 外側の半径 $R$、軸の半径 $r$ の糸巻きの、軸の下側から出た糸を、水平から角 $\phi$ の向きに引きます。糸巻きが引く向きに転がるための $\phi$ の条件を求めましょう。また、$R=2r$ のとき、その境の角は何度ですか。
\bkpans{床にふれた点のまわりのトルク $F(R\cos\phi-r)$ が正のとき。$\cos\phi>\dfrac{r}{R}$。$R=2r$ なら境は\ang{60}（\ang{60}より小さい角で引く向きへ転がる）}
\bkproblem 回転台の中心から\qty{0.40}{m}の所に、はば\qty{0.10}{m}、高さ\qty{0.30}{m}の一様な直方体の箱を立てました（はばは半径の向き）。静止摩擦係数を0.50とします。台の回転を少しずつ速くすると、箱はすべり出すのと倒れるのと、どちらが先ですか。そのときの角速度を求めましょう（箱の大きさは半径に比べて小さいとして、遠心力は重心にかかるとする）。
\bkpans{すべる $\omega=\sqrt{\dfrac{0.50\times9.8}{0.40}}=\qty{3.5}{rad/s}$、倒れる $mr\omega^2\times0.15=mg\times0.05$ より $\omega=\sqrt{\dfrac{9.8\times0.05}{0.40\times0.15}}\approx\qty{2.9}{rad/s}$。倒れるほうが先（約\qty{2.9}{rad/s}）}

\begin{furikaeri}
\begin{checklist}
\item 転がる車輪の点の速さを「中心の速さ＋回る速さ」で求め、$v=r\omega$ を使える。\bkfurref{ph:力-7:sec1}
\item 2本の腕の関節の角から手の位置を、手の位置から関節の角を求められる。歯車の速さとトルクを求められる。\bkfurref{ph:力-7:sec3}
\item 平面運動の式と回る座標の式から、斜面を転がる加速度や、遠心力・コリオリの力を導ける。\bkfurref{ph:力-7:sec5b}
\end{checklist}
\end{furikaeri}
\bkfukidashi{sora}{ロボットの腕も、自転車の車輪も、地球の上の風も、「回る」と「進む」を分けて考えれば、同じ道具で計算できるんだね。}
